% An area law for four-dimensional SU(2) lattice gauge theory up to beta ~ 0.813.
% Version 1.2, 12 November 2026 (literature situating only; no theorem, constant or certificate changed). The previous version is v1.1, doi:10.5281/zenodo.23877371. Compile: pdflatex, three passes.
\documentclass[11pt,a4paper]{article}
\usepackage[T1]{fontenc}
\usepackage[utf8]{inputenc}
\usepackage{lmodern}
\usepackage[english]{babel}
\usepackage{amsmath,amssymb,amsthm}
\usepackage{booktabs,array}
\usepackage{geometry}
\geometry{margin=2.6cm}
\usepackage[colorlinks=true,linkcolor=blue,citecolor=blue,urlcolor=blue]{hyperref}
\usepackage{xurl}% \path{} for hashes and file names, breakable anywhere
\setlength{\emergencystretch}{2em}% v1.1: paragraphs with long \texttt names break without overfull lines

\hypersetup{pdftitle={An area law for four-dimensional SU(2) lattice gauge theory up to beta = 0.813},pdfauthor={K\'evin R\'emondi\`ere}}

\theoremstyle{plain}
\newtheorem{theorem}{Theorem}[section]
\newtheorem{proposition}[theorem]{Proposition}
\newtheorem{lemma}[theorem]{Lemma}
\newtheorem{corollary}[theorem]{Corollary}
\theoremstyle{remark}
\newtheorem{remark}[theorem]{Remark}
\numberwithin{equation}{section}

\newcommand{\R}{\mathbb{R}}
\newcommand{\Z}{\mathbb{Z}}
\newcommand{\N}{\mathbb{N}}
\newcommand{\E}{\mathbb{E}}
\newcommand{\Tr}{\operatorname{Tr}}
\newcommand{\ReTr}{\operatorname{Re}\operatorname{Tr}}
\newcommand{\cT}{\mathcal{T}}
\newcommand{\cH}{\mathcal{H}}
\newcommand{\At}{\widetilde A}
\newcommand{\rt}{\widetilde r}
\newcommand{\rb}{\bar r}
\newcommand{\Gt}{\widetilde G}
\newcommand{\st}[1]{\textsf{[#1]}}
\newcommand{\bs}{\boldsymbol}

\title{An area law for four-dimensional $SU(2)$ lattice gauge theory\\ up to $\beta\approx0.813$}
\author{K\'evin R\'emondi\`ere\thanks{Independent researcher, Oloron-Sainte-Marie, France. ORCID
\href{https://orcid.org/0009-0008-2443-7166}{0009-0008-2443-7166}. The proof, the programs and this note were prepared with
the assistance of AI models (Claude Opus~5.5, Anthropic; Kimi K3.1 Pro, Moonshot AI; DeepSeek V4 Pro, DeepSeek); their roles are
described in the statement on the use of AI at the end of the note. For version~1.2, the passages on the literature were revised on 11 November 2026 with the assistance of GPT-6.1 Sol (OpenAI), after a review of the draft of 10 November 2026 by Qwen3.8 (Qwen team, run locally); the final text was then read against the sources by Claude (Anthropic) agents on 11 November 2026; no statement, proof, constant or certificate was modified since version~1.1.}}
\date{Version 1.2, 12 November 2026 (version 1.1: 2 November 2026)}

\begin{document}
\maketitle

\begin{abstract}
We consider the $SU(2)$ Wilson lattice gauge theory on the periodic torus $(\Z/L\Z)^4$, with action
$\beta\sum_p(1-\tfrac12\ReTr U_p)$. We prove that for $0\le\beta\le0.81348$ every planar rectangular Wilson loop with sides
$R\le L/2$ and $T\le L-1$ satisfies $|\langle W\rangle|\le C^T\theta^{RT}$, with explicit constants $\theta<1$ and $C$ independent
of $L$; in particular the string tension is positive. The proof follows the method of Borgs and Seiler (1983, Appendix~B): the slab
decomposition of Durhuus and Fr\"ohlich, a random-walk expansion of the slab $\sigma$-models and the positivity of the transfer
matrix, which yields the range $\beta<12/17\approx0.7059$. We refine it in two ways. (I)~The ratio of slab partition
functions is controlled by a strict operator inequality $\cT^{(\lambda)}\le c(\lambda)\,\cT$, where $c(\lambda)<1$ comes from the
Peter--Weyl spectrum (the spin-$0$ mode is extremal), followed by a site-by-site Cauchy--Schwarz factorisation. (II)~The resulting
weighted walks are counted with a memory of $11$ steps, through a transfer matrix on $7\,821\,056$ symmetry classes whose spectral
radius is bounded by an integer Perron certificate. This refines the threshold of the Borgs--Seiler method from $12/17$ to
$0.81348$ ($+15.2\%$), with bounds for each rectangle on tori. Every numerical inequality of the proof is verified in exact integer or rational
arithmetic, and the certificate is re-verified by two further programs, one of them written separately for an adversarial
verification. With memory $13$ the same argument gives
$\beta\le0.81498$. The result is computer-assisted, holds on tori only, and has been checked by machine (independent programs,
reviews by language models, and a partial formalisation in Lean that assumes the analytic reduction), not by a human referee. The value $0.81348$ lies between the estimate $4/5$ announced without proof by
Fr\"ohlich (1980), in the reading discussed in \S\ref{sec:heur}, and the value $0.8222$ given by the non-backtracking heuristic
$5\,I_2(\beta)/I_1(\beta)<1$, which we evaluate in \S\ref{sec:heur}. Version~1.1 adds a tilted form of the walk bound
(Theorem~\ref{thm:tilted}): with the same threshold and five further integer certificates, the certified lower bound on the string
tension at $\beta=0.81348$ rises from $7.05\cdot10^{-6}$ to $0.004989$, and to $0.2253$ for $\beta\le0.80$, with explicit constants on
every torus.
\end{abstract}

\paragraph{Status labels.}
\st{proved}: complete argument written in this note, or in the cited source at the location given;
\st{certified}: inequality between explicit numbers, verified by a program in exact integer or rational arithmetic and re-verified by
at least one other program (\S\ref{sec:cap});
\st{numerical}: floating-point or extrapolated value, never used in a proof;
\st{read}: statement read in the cited source, location given.

\tableofcontents

\section{Introduction}\label{sec:intro}

\subsection{Context}
Wilson's criterion \cite{Wil74} expresses the confinement of static quarks by the area law for Wilson loops. For lattice gauge
theories with a compact gauge group, the area law at strong coupling was proved by Osterwalder and Seiler \cite[Thm.~5.1]{OS78} by
a cluster expansion, in a range of small $\beta$ that is not made explicit. For the four-dimensional $SU(2)$ theory, explicit ranges
come from the decomposition of Durhuus and Fr\"ohlich \cite{DF80}: conditionally on the horizontal links, the vertical links of each
time slab form a three-dimensional $O(4)$ $\sigma$-model with random twists, and a rectangular Wilson loop factorises over the slabs.
Table~\ref{tab:ranges} lists the explicit ranges we know, in the normalisation \eqref{eq:wilson} below.

\begin{table}[ht]
\centering\footnotesize
\begin{tabular}{@{}>{\raggedright\arraybackslash}p{0.15\textwidth}>{\raggedright\arraybackslash}p{0.41\textwidth}%
>{\raggedright\arraybackslash}p{0.18\textwidth}>{\raggedright\arraybackslash}p{0.15\textwidth}@{}}
\toprule
range of $\beta$ & setting & source & status \\
\midrule
small, not explicit & planar loops; cluster expansion & \cite[Thm.~5.1]{OS78} & proved \\
$\lesssim1/8$ & refined cluster expansion of Guerra et al., as quoted in \cite{BS83} & \cite[p.~343]{BS83} & quoted \\
$<1/6$ & tori, uniformly in $L$ & \cite[Thm.~1.6]{CNS25} & proved \\
$<1/6$ & written argument of \cite{DF80}, after correction of a factor $2$\,$^{a}$ & \cite{DF80}, \cite[\S7.2]{D7} & proved \\
$\approx0.2217$ & comparison with $\Z_2$ and Ising slabs \cite{Fro79}\,$^{b}$ & \cite[\S7.2]{D7} & not checked \\
$<1/\sqrt3\approx0.577$ & tori; slab $\sigma$-model, geodesic one-link constant & \cite[Part~I, Cor.~6.5]{Raj26a} & \st{read}; proof not checked here \\
$<2/3$ & tori, free boundary, every Gibbs measure; one-site Dobrushin & \cite[Thm.~7.1]{D7} & proved \\
$<2/3$ & tori; slab model with the chordal one-link constant $\frac14$ & \cite[Part~II, Cor.~7.7(iii)]{Raj26a} & \st{read}; proof not checked here \\
$<12/17\approx0.7059$ & Polyakov loops at nonzero temperature, uniformly in $L_0,J_M$\,$^{c}$ & \cite[(II.55), p.~343; App.~B, pp.~373--377]{BS83} & proved \\
$\approx0.730$\,$^{e}$ & leading-order window of a decimated slab $\sigma$-model, ``$\sqrt{16/30}$'' & \cite[\S8, item (iv)]{Raj26b} & \st{read}; heuristic \\
$<4/5$\,$^{d}$ & averaged (``in mean'') estimate, ``only a rough estimate'' & \cite{Fro80a,Fro80b} & heuristic \\
$\le0.81348$ & tori, uniformly in $L$; memory $11$ & Theorem~\ref{thm:main} & proved, computer-assisted \\
$\le0.81498$ & same; memory $13$ & Remark~\ref{rem:m13} & proved, computer-assisted \\
$<0.8222$ & non-backtracking heuristic $5u(\beta)<1$ & \S\ref{sec:heur} & heuristic \\
all $\beta$ (claim) & preprint; gaps discussed in \cite{IS08} & \cite{Tom07} & contested \\
\bottomrule
\end{tabular}
\caption{Explicit ranges of the area law for four-dimensional $SU(2)$ in the Wilson normalisation $\beta=4/g^2$.
$^{a}$~The literal statement of \cite{DF80} is $\beta<1/3$, and $\beta<2/3$ is announced there without proof.
$^{b}$~Combined with the sharpness of the three-dimensional Ising transition and $K_c^{\mathrm{Ising}}(\Z^3)=0.221654626(5)$
\cite{FXL18}, as described in \cite[\S7.2]{D7}; we have not seen it written out. $^{c}$~See Remark~\ref{rem:bs-statement}; the same method gives rectangular Wilson loops at zero temperature (\S\ref{sec:bs}).
$^{d}$~In the reading of \S\ref{sec:heur}; the literal reading gives $8/5$.
$^{e}$~This heuristic window concerns a decimated slab model; the numerical value $\approx0.730$ for improvement~(II) alone in \S1.2 concerns our finite-memory walk count. No identification between the two estimates is asserted.}
\label{tab:ranges}
\end{table}

Two observations frame this note. First, an argument that proves exponential decay of the slab correlations \emph{uniformly in the
twists} cannot go beyond the critical coupling of the untwisted $O(4)$ model on $\Z^3$, $\beta_c\approx0.936$ according to
Monte Carlo estimates \cite{KK95,CH98}, since the twists may be trivial; the one-site Dobrushin criterion stops at $2/3$ exactly \cite[Remark~7.6]{D7}. In the same spirit, Proposition~7.1 of the second preprint of the record \cite{Raj26b} states a cap $\beta\le4G_3(0)$ for properties of the ungauged Wilson specification that imply uniqueness of the infinite-volume Gibbs measure and are inherited by fixing links. Its examples are the single-site Dobrushin condition and the strong-mixing form of complete analyticity, uniformly in all finite volumes and boundary conditions; its proof applies the infrared bound of Fr\"ohlich, Simon and Spencer to the three-dimensional $O(4)$ model of the time-like links (\st{read}).
Second, Borgs and Seiler \cite[Appendix~B]{BS83} do not need uniformity in the twists: after a random-walk expansion of the slab
correlations, the average over the horizontal links is taken inside a trace of products of transfer matrices, and only the positivity
of the transfer matrix and a bound on a ratio of slab partition functions are used. Their bound on that ratio is $1$ (monotonicity in
the coupling), and their count of walks only distinguishes an immediate return; the result is $J_E<(d-\frac16)^{-1}$, that is
$\beta<12/17$ for $d=3$ space dimensions \cite[(II.55), p.~343, and p.~376]{BS83} \st{read}. They state it for correlations of Polyakov loops at
nonzero temperature (Remark~\ref{rem:bs-statement}); the argument applies to rectangular Wilson loops at zero temperature
(\S\ref{sec:bs}), and we call $\beta<12/17$ the \emph{threshold of the Borgs--Seiler method}.

\subsection{Result}
Theorem~\ref{thm:main} proves the area law on the tori $(\Z/L\Z)^4$, $L\ge3$, uniformly in $L$, for $0\le\beta\le0.81348$;
Remark~\ref{rem:m13} records $\beta\le0.81498$ with a larger certificate. Relative to $12/17\approx0.70588$ the gain is $15.2\%$
($15.5\%$). It comes from two refinements of the Borgs--Seiler argument, which can be used separately or together.
\begin{itemize}
\item[(I)] \emph{A strict operator inequality} (\S\ref{sec:I}). The transfer matrix $\cT^{(\lambda)}$ of a slab with the reduced
couplings produced by the walk expansion satisfies $\cT^{(\lambda)}\le c(\lambda)\,\cT$, where $c(\lambda)$ is a product over links
of ratios of the spin-$0$ eigenvalues of the one-link convolution operators; for $J>0$, $c(\lambda)<1$ as soon as
$\lambda\not\equiv0$. Borgs
and Seiler use $c\le1$. A Cauchy--Schwarz inequality turns $c(\lambda)$ into a product of factors attached to the sites of the walk,
which lowers the weight of every visit.
\item[(II)] \emph{Walks with memory} (\S\ref{sec:II}). The weighted walks are counted with the visits seen in a window of the last
$m=11$ steps, through a transfer matrix on $6^{11}$ step sequences reduced by the $48$ symmetries of $\Z^3$; its spectral radius is
bounded by an integer Perron vector, checked in $128$-bit integer arithmetic.
\end{itemize}
Separately, (I) gives about $0.782$ and (II) about $0.730$ \st{numerical}; together they give $0.81348$. A monotonicity property in
the coupling (Lemma~\ref{lem:monotone}) lets a single certificate cover the whole interval $[0,0.81348]$.

Version~1.1 adds a tilted form of the walk bound (\S\ref{sec:tilted}). Weighting each step of the walks by
$e^{\eta\langle\hat e,d\rangle}$, where $\hat e$ is the direction of the side of length $R$, and using the parity of walks in $\Z^3$,
Theorem~\ref{thm:tilted} keeps the threshold and, with five further memory-$11$ certificates, bounds the string tension below by
$0.004989$ at $\beta_*$ instead of $7.05\cdot10^{-6}$, and by $0.2253$ for $\beta\le0.80$ (Table~\ref{tab:tilted}); at $\beta_*$ the bound
becomes non-trivial from $R\approx3.1\cdot10^{3}$ on every torus ($2.8\cdot10^{3}$ in the weak limits) instead of $1.9\cdot10^{6}$.

\subsection{What is new and what is not}
To our knowledge, the widest explicit confinement range already published for the $SU(2)$ Wilson theory in four dimensions
is $\beta<12/17\approx0.7059$, stated by Borgs and Seiler \cite[(II.55), p.~343]{BS83} and proved in their Appendix~B (pp.~373--377).
It concerns correlations of Polyakov loops at every temperature and every magnetic coupling; Wilson's string tension is controlled through \cite[(II.50), p.~341]{BS83}, and the
authors do not state a bound for each rectangular loop. Theorem~\ref{thm:main} is a refinement of the threshold of the Borgs--Seiler
method, from $12/17$ to $0.81348$ ($+15.2\%$), with bounds for each rectangle on tori, uniform in the volume. It is not an independent
method: the new ingredients are the strict inequality (I) with its site-by-site factorisation, its combination with finite-memory
counting, and the certificate.

Among the other public statements, Part~I of Rajabu \cite[Cor.~6.5]{Raj26a} states the area law on tori for $\beta<1/\sqrt3\approx0.577$ and Part~II \cite[Cor.~7.7(iii)]{Raj26a} for $\beta<2/3$, using the slab $\sigma$-model of \cite{DF80} and a Dobrushin contraction with the geodesic, resp.\ the chordal one-link constant. The Zenodo record was created on 27 October 2026, 13:09:49 UTC; we have not checked these proofs. Version~2.2 of his preprint \cite[\S8, item (iv)]{Raj26b} mentions, as a direction and not as a theorem, a leading-order window of about $0.730$ for a decimated slab $\sigma$-model. Its record has \texttt{publication\_date} 28 October 2026 and was created on 29 October 2026, 00:48 UTC. The range $2/3$ of \cite[Thm.~7.1]{D7} was published on 1 November 2026, after Rajabu's record of 27 October. No priority claim is made for any threshold. In this note's normalisation, $0.81348$ exceeds $2/3$ by $22.0\%$ and the threshold $12/17$ of \cite{BS83} by $15.2\%$; these are comparisons of ranges, with the settings specified in Table~\ref{tab:ranges}.

Not new: the slab decomposition \cite{DF80}; the random-walk expansion of the slabs and the use of transfer-matrix positivity
\cite{BS83}, which we follow closely and write out with the details we need (\S\ref{sec:bs}); and finite-state transfer-matrix bounds for the connective constant of
self-avoiding walks \cite{PT00}. For related upper bounds, see also
\cite{Alm93}. Improvement~(I) is natural (it only uses the Peter--Weyl decomposition of the one-link
convolution operators), and we cannot exclude that it appears in work we have not seen. Borgs and Seiler themselves point out that
their bound could be sharpened by the perimeter decay of the characters, without giving a number \cite[p.~377]{BS83}; we do not use
this route. Tomboulis \cite{Tom07} claims confinement in four-dimensional $SU(2)$ lattice gauge theory for all values of the
coupling; Ito and Seiler \cite{IS08} discuss gaps and open problems in its argument. We do
not count it as an established range. The number itself is close to the
heuristic values discussed in \S\ref{sec:heur}: it exceeds the estimate $4/5$ announced without proof by Fr\"ohlich
\cite{Fro80a,Fro80b} in one reading of his normalisation, and it is below the value $0.8222$ obtained here from the non-backtracking
heuristic, of which $4/5$ is the small-$\beta$ linearisation. Like every explicit strong-coupling range, the result says nothing about the weak-coupling regime
relevant to the continuum limit.

\subsection{Status of the proof}\label{sec:status}
The analytic part (\S\S\ref{sec:bs}--\ref{sec:II}) is proved on paper. Its numerical inputs, namely rational enclosures of twelve
one-dimensional integrals and the inequality $T_{11}H\le\rho'H$ on $7\,821\,056$ symmetry classes, are verified in exact rational and
integer arithmetic and re-verified by other programs, written with independent implementation choices or in separate sessions
(\S\ref{sec:cap}). The proof and the certificates have been
checked by machine only: by independent programs, by adversarial reviews carried out by language models (see the statement on the
use of AI) and, in part, in Lean; they have not been read by a human referee. The Lean formalisation (\S\ref{sec:repro}) covers the
combinatorial part, the certificate checker and the final arithmetic; the analytic reduction of \S\S\ref{sec:bs}--\ref{sec:I}, the
monotonicity of Lemma~\ref{lem:monotone} and the enclosures of the weights are assumed there. Theorem~\ref{thm:tilted}
(version~1.1) uses the same analytic part and five further integer certificates, each checked by two programs
(\S\ref{sec:tiltcap}); it is not formalised in Lean.

\section{Setting and statement}\label{sec:statement}

\subsection{Model and normalisation}
Let $L\ge3$ and $\Lambda_L:=(\Z/L\Z)^4$. Each link $(z,\mu)$, $z\in\Lambda_L$, $\mu\in\{1,2,3,4\}$, carries $U_{z,\mu}\in SU(2)$.
A plaquette is $p=(z;\mu,\nu)$ with $\mu<\nu$, each geometric plaquette being counted once, and
$U_p:=U_{z,\mu}U_{z+e_\mu,\nu}U_{z+e_\nu,\mu}^{-1}U_{z,\nu}^{-1}$. For $\beta\ge0$ the Wilson measure is
\begin{equation}\label{eq:wilson}
\mu_{L,\beta}(dU)=\frac1Z\exp\Big(\beta\sum_p\tfrac12\ReTr U_p\Big)\prod_\ell dU_\ell ,
\end{equation}
with normalised Haar measure $dU_\ell$; up to a constant factor the weight is $\exp(-\beta\sum_p(1-\frac12\ReTr U_p))$. This is the
Wilson normalisation, $\beta=4/g^2$. We write $\langle\cdot\rangle_{L,\beta}$ for the expectation and
\begin{equation}\label{eq:J}
J:=\beta/2
\end{equation}
for the coefficient of $\ReTr U_p$ in the exponent of \eqref{eq:wilson}. In other normalisations: the electric coupling $J_E$ of
Borgs and Seiler \cite[(II.2)--(II.4), p.~332]{BS83} at isotropic lattice spacing is $J_E=2/g^2=J$ \st{read}; the coupling of
\cite[(1.13)--(1.16)]{DF80} and of \cite{Fro79} (coefficient of $\operatorname{Re}\chi$, $\chi=\Tr$) is $J$; the 't~Hooft coupling of
\cite{CNS25} is $\beta/4$.

For a closed lattice path $\gamma$, $U_\gamma$ is the ordered product of its link variables and $W_\gamma:=\frac12\Tr U_\gamma\in
[-1,1]$. A \emph{rectangular loop with sides $R$ along $e_\mu$ and $T$ along $e_\nu$} ($\mu\ne\nu$, $R,T\ge1$) is the boundary of a
planar $R\times T$ rectangle of plaquettes.

\subsection{Main result}

\begin{theorem}[computer-assisted]\label{thm:main}
Let
\[
J_*:=\frac{40674}{100000},\qquad \beta_*:=2J_*=0.81348,\qquad \rho'_*:=\frac{2639854083}{2^{30}},
\]
\[
\theta_*:=J_*\rho'_*=\frac{53686712485971}{53687091200000}=1-7.0540\ldots\times10^{-6},
\]
and let $C_*:=K_0/(1-\theta_*)$, with $K_0$ defined in \eqref{eq:K0}; numerically $K_0=3.97805\ldots$ and
$C_*=5.6393\ldots\times10^{5}$. For every $\beta\in[0,\beta_*]$, every $L\ge3$ and every rectangular loop $\gamma$ with sides $R$
along $e_\mu$ and $T$ along $e_\nu$ such that $1\le R\le L/2$ and $1\le T\le L-1$,
\begin{equation}\label{eq:main}
|\langle W_\gamma\rangle_{L,\beta}|\ \le\ C_*^{\,T}\,\theta_*^{\,RT}.
\end{equation}
\end{theorem}
\noindent Status: \st{proved} (\S\S\ref{sec:bs}--\ref{sec:II}), with two \st{certified} numerical inputs
(Proposition~\ref{prop:cert}). The constants depend neither on $L$ nor on $\beta\in[0,\beta_*]$. By the symmetries of the lattice the
two sides may be exchanged: \eqref{eq:main} with $R$ and $T$ exchanged holds when $T\le L/2$ and $R\le L-1$.

\begin{corollary}[string tension]\label{cor:tension}
Let $0<\beta\le\beta_*$. For all $R,T\ge1$ and $L\ge\max(2R,T+1)$,
\begin{equation}\label{eq:sigma}
-\frac1{RT}\log|\langle W_{R\times T}\rangle_{L,\beta}|\ \ge\ -\log\theta_*-\frac{\log C_*}{R}.
\end{equation}
Hence
\[
\sigma(\beta):=\liminf_{R\to\infty}\ \inf_{T\ge1}\ \inf_{L\ge\max(2R,T+1)}\Big(-\frac1{RT}\log|\langle W_{R\times T}\rangle_{L,\beta}|
\Big)\ \ge\ -\log\theta_*=7.05\ldots\times10^{-6}>0,
\]
and any definition of the string tension by a limit in which $R\to\infty$ (with $T\to\infty$
or not, for instance $\liminf_{R,T\to\infty}$) satisfies the same bound. The bound \eqref{eq:main} also holds for every weak limit point of $(\mu_{L,\beta})_{L}$ as
$L\to\infty$.
\end{corollary}
\begin{proof}
\eqref{eq:sigma} is \eqref{eq:main} rewritten. The limit $R\to\infty$ is needed to remove the term $(\log C_*)/R$. For a fixed loop,
$W_\gamma$ is a bounded continuous function of finitely many links, so \eqref{eq:main} passes to weak limits.
\end{proof}

\begin{remark}[memory $13$]\label{rem:m13}
The same proof with memory $m=13$ gives \eqref{eq:main} for $\beta\le\beta_*^{(13)}:=0.81498$, with
\[
J_*^{(13)}=\frac{40749}{100000},\qquad \rho'^{(13)}_*=\frac{2635005744}{2^{30}},\qquad
\theta_*^{(13)}=\frac{6710865566391}{6710886400000}=1-3.1044\ldots\times10^{-6},
\]
$K_0^{(13)}=5.2781\ldots$ and $C_*^{(13)}=1.7002\ldots\times10^{6}$ \st{certified}. Its certificate has $276\,294\,656$ symmetry
classes and $4.4$\,GB; it is identified by its SHA-256 in \S\ref{sec:repro} and can be regenerated. We state this improvement as a
remark because the certificate is not part of the supplementary material.
\end{remark}

\begin{remark}[string-tension bounds]\label{rem:table}
The same argument at $J=\beta_i/2$, with a certificate of memory $11$ for each $\beta_i$, gives \eqref{eq:main} for all
$\beta\le\beta_i$ with the constants of Table~\ref{tab:tension} \st{certified}. For $\beta\le0.66$ the bound
$\sigma\ge\operatorname{arcosh}(2/\beta-2)$ of \cite[Cor.~7.2]{D7} is better (for instance $0.795$ against $0.256$ at $\beta=0.6$); it
tends to $0$ as $\beta\uparrow2/3$. Corollary~\ref{cor:tilted} (version~1.1) gives $\sigma\ge0.6080$ for all $\beta\le0.70588$, which is better than
$\operatorname{arcosh}(2/\beta-2)$ for $\beta>0.627$. Without the tilt, the walk method is useful for the range rather than for the size of
the tension.
Near the threshold the bound \eqref{eq:main} is non-trivial only at large distances: $C^T\theta^{RT}<1$ if and only if
$R>R_{\min}:=\log C/(-\log\theta)$. From the constants of Table~\ref{tab:tension}, $R_{\min}\approx12$ for $\beta\le0.60$,
$143$ for $\beta\le0.78$, $430$ for $\beta\le0.80$, $2.1\cdot10^{3}$ for $\beta\le0.81$ and $1.9\cdot10^{6}$ for $\beta\le\beta_*$
(with $C_*=5.6\cdot10^5$ and $-\log\theta_*=7.1\cdot10^{-6}$); at $\beta_*$ the bound is purely asymptotic. With Theorem~\ref{thm:tilted} (version~1.1), Table~\ref{tab:tilted} gives instead
$R_{\min}\le8.5$, $36.5$, $85.5$ and $3\,107$ on every torus ($6.3$, $30.3$, $72.4$ and $2\,829$ in the weak limits) for $\beta\le0.70588$,
$0.80$, $0.81$ and $\beta_*$.
\end{remark}

\begin{table}[ht]
\centering\small
\begin{tabular}{@{}llll@{}}
\toprule
$\beta_i$ & $\theta_i$ & $-\log\theta_i$ & $C_i$\\
\midrule
$0.60$ & $0.77398$ & $0.2562$ & $19.3$\\
$0.66$ & $0.84102$ & $0.1731$ & $26.8$\\
$0.70588$ & $0.89048$ & $0.1160$ & $38.2$\\
$0.72$ & $0.90538$ & $0.0994$ & $43.9$\\
$0.74$ & $0.92621$ & $0.0767$ & $55.8$\\
$0.76$ & $0.94672$ & $0.0547$ & $76.6$\\
$0.78$ & $0.96692$ & $0.0336$ & $122$\\
$0.80$ & $0.98679$ & $0.0133$ & $303$\\
$0.81$ & $0.99660$ & $0.0034$ & $1.2\cdot10^{3}$\\
$0.81348$ & $0.9999929$ & $7.1\cdot10^{-6}$ & $5.6\cdot10^{5}$\\
\bottomrule
\end{tabular}
\caption{Constants of \eqref{eq:main} valid for all $\beta\le\beta_i$ (memory $11$). Each row rests on its own weights and integer
certificate, produced and checked as in \S\ref{sec:cap} and re-verified on all $6^{11}$ states by an independent program; only the
last one is in the supplementary material, the others can be regenerated with the programs of \S\ref{sec:repro}. The weights files
of the rows $0.70588$, $0.80$ and $0.81$ are in \texttt{supplement\_v1.1} (version~1.1).}
\label{tab:tension}
\end{table}

\section{The method of Borgs and Seiler}\label{sec:bs}

This section writes out, in our notation and for rectangular Wilson loops on tori, the argument of \cite[Appendix~B]{BS83}. We take
$\nu=4$ (the side of length $T$ is vertical) and $\mu=1$; the general case follows by symmetry.

\subsection{Slabs}
Let $\iota:S^3\to SU(2)$, $\iota(q)=\begin{pmatrix}q^0+iq^3&-q^1+iq^2\\ q^1+iq^2&q^0-iq^3\end{pmatrix}$ (the parametrisation of
\cite[(5.1)]{DF80}); it maps the normalised surface measure $\bar\sigma$ of $S^3\subset\R^4$ onto Haar measure, and
$\frac12\ReTr(\iota(q)^{-1}\iota(r))=\langle q,r\rangle$. For $A,B\in SU(2)$ the map $r\mapsto\iota^{-1}(A\,\iota(r)B^{-1})$ is the
restriction to $S^3$ of a matrix $O_{A,B}\in SO(4)$, with $O_{A,B}O_{A',B'}=O_{AA',BB'}$ and $O_{A,B}^{\mathsf T}=O_{A^{-1},B^{-1}}$.

Let $X:=(\Z/L\Z)^3$, with its set $E_X$ of $3L^3$ nearest-neighbour edges; for $L\ge3$ it is a simple $6$-regular graph. The
\emph{vertical} links are $Q^k_x:=U_{(x,k),4}$, $x\in X$, $k\in\Z/L\Z$, and $q^k_x:=\iota^{-1}(Q^k_x)$; the other links are
\emph{horizontal}, and $\cH$ is the $\sigma$-algebra they generate. A vertical plaquette $p=((x,k);\mu,4)$, $\mu\le3$, has
$U_p=A\,Q^k_{x+e_\mu}B^{-1}(Q^k_x)^{-1}$ with $A=U_{(x,k),\mu}$ (layer $k$) and $B=U_{(x,k+1),\mu}$ (layer $k+1$), so that
$\frac12\ReTr U_p=\langle q^k_x,O_{A,B}\,q^k_{x+e_\mu}\rangle$.

\begin{lemma}[slabs; {\cite[(1.17)--(1.19)]{DF80}}, {\cite[Lemma~7.4]{D7}}]\label{lem:slab}
Conditionally on $\cH$, the families $(q^k_x)_{x\in X}$, $k\in\Z/L\Z$, are independent, and $(q^k_x)_x$ has the law
\begin{equation}\label{eq:slab}
\mu_k(dq)=\frac1{Z_k}\exp\Big(2J\sum_{x\in X,\,\mu\le3}\langle q_x,O^k_{x,x+e_\mu}\,q_{x+e_\mu}\rangle\Big)\prod_{x\in X}\bar\sigma(dq_x),
\end{equation}
with $O^k_{x,x+e_\mu}:=O_{U_{(x,k),\mu},\,U_{(x,k+1),\mu}}$.
\end{lemma}
\begin{proof}
A horizontal plaquette contains no vertical link, and a vertical plaquette contains two vertical links of the same slab; the density
\eqref{eq:wilson} therefore factorises over the slabs given $\cH$, and $\beta\cdot\frac12\ReTr U_p=2J\langle q_x,O\,q_{x+e_\mu}\rangle$.
\end{proof}
\noindent We write $O_{yx}:=O_{xy}^{\mathsf T}$ for the reversed edge, and $Z_k(\bs J)$ for the normalisation of \eqref{eq:slab} with
the coupling $J$ of each edge $\ell$ replaced by $J_\ell$; it is a function of the horizontal links of layers $k$ and $k+1$.

\subsection{Random-walk expansion of a slab}
\begin{lemma}[integration by parts on a sphere]\label{lem:ipp}
For $F\in C^1(\R^4)$, $c>0$ and $a\in\{0,1,2,3\}$,
$\displaystyle\int s_aF(s)\,\delta(|s|^2-c)\,d^4s=\frac12\int_0^c dc'\int\partial_aF(s)\,\delta(|s|^2-c')\,d^4s$.
\end{lemma}
\begin{proof}
Let $\Theta(c):=\int_{|s|^2\le c}\partial_aF\,d^4s$. By the divergence theorem $\Theta(c)=\int_{|s|=\sqrt c}F\,s_a/\sqrt c\,dS$; by the
coarea formula $\int g\,\delta(|s|^2-c)\,d^4s=(2\sqrt c)^{-1}\int_{|s|=\sqrt c}g\,dS$ and $\Theta'(c)=\int\partial_aF\,\delta(|s|^2-c)\,
d^4s$. Hence the left side equals $\Theta(c)/2=\frac12\int_0^c\Theta'$.
\end{proof}
\noindent The factor $\frac12$ does not appear in the printed form of \cite[Lemma~B.1(i), p.~374]{BS83}; it is compensated elsewhere in their
computation, and the net weight per step is $J$ in both (an independent check of the normalisation is in \S\ref{sec:checks}).

Let $G$ be a finite simple graph, $J\ge0$, $O_{xy}\in SO(4)$ for each ordered pair of neighbours with $O_{yx}=O_{xy}^{\mathsf T}$, and
\[
Z(\bs J):=\int\exp\Big(2\sum_{\{u,v\}\in E_G}J_{uv}\langle q_u,O_{uv}q_v\rangle\Big)\prod_{z\in G}\bar\sigma(dq_z)
\qquad(\bs J=(J_{uv})_{\{u,v\}\in E_G}),
\]
and let $\mu$ be the corresponding probability measure for the constant couplings $J_{uv}\equiv J$.
A \emph{walk} $\omega:x\to y$ is a sequence $\omega_0=x,\omega_1,\dots,\omega_{|\omega|}=y$ of sites with $\omega_{i+1}$ a neighbour
of $\omega_i$; returns and self-intersections are allowed. Put $O_\omega:=O_{\omega_0\omega_1}\cdots O_{\omega_{|\omega|-1}
\omega_{|\omega|}}$, $n_z(\omega):=\#\{i:\omega_i=z\}$, $V(\omega):=\{z:n_z(\omega)\ge1\}$,
\begin{equation}\label{eq:fn}
f_n(\lambda):=\frac{\lambda^{n-1}(1-\lambda)}{(n-1)!}\quad(n\ge1,\ \lambda\in[0,1]),\qquad F_\omega(\lambda):=\prod_{z\in V(\omega)}
f_{n_z(\omega)}(\lambda_z),
\end{equation}
and, for $\lambda\in[0,1]^{V(\omega)}$ extended by $\lambda_u:=0$ for $u\notin V(\omega)$, the reduced couplings
$J^\lambda_{uv}:=J\sqrt{(1-\lambda_u)(1-\lambda_v)}$.

\begin{proposition}[walk representation; {\cite[(B.1)--(B.8), pp.~374--375]{BS83}}]\label{prop:walk}
For $x\ne y$ in $G$ and $a,b\in\{0,\dots,3\}$,
\begin{equation}\label{eq:walk}
\E_\mu[q_{x,a}\,q_{y,b}]=\frac12\sum_{\omega:x\to y}J^{|\omega|}\,(O_\omega)_{ab}\int_{[0,1]^{V(\omega)}}F_\omega(\lambda)\,
\frac{Z(\bs J^\lambda)}{Z(J)}\,d\lambda ,
\end{equation}
and the series converges absolutely.
\end{proposition}
\begin{proof}
Write $\bar\sigma(dq)\propto\delta(|s|^2-1)\,d^4s$ (the constant cancels in the ratio) and $e^{H}$ for the weight, so that
$\partial_{s_z,c}H=2J\sum_{u\sim z}(O_{zu}s_u)_c$.
(i) Apply Lemma~\ref{lem:ipp} to the factor $s_{x,a}$: since $\frac12\,\partial_{x,a}e^H=J\sum_{u\sim x}(O_{xu}s_u)_ae^H$, this
produces a step to a neighbour $u$ with weight $\frac12\cdot2J=J$, a new linear factor $(O_{xu}s_u)_a$, and replaces the constraint
at $x$ by $\delta(|s_x|^2-(1-\lambda_x))$ with $\lambda_x\in[0,1]$.
(ii) Repeat with the linear factor at the current site. When the current site is $y$, the derivative may also hit $s_{y,b}$; this
gives $\frac12\delta_{cb}$ and ends the walk, which is the global factor $\frac12$ of \eqref{eq:walk}. The derivative of $H$ always
carries the unreduced coupling $J$.
(iii) At the $n$-th visit of a site $z$ the new parameter runs over $[\lambda_z^{(n-1)},1]$ and the integrand depends only on the last
one; the ordered simplex gives $\lambda^{n-1}/(n-1)!$.
(iv) At the end, the change of variables $s_z=\sqrt{1-\lambda_z}\,s'_z$ gives $\delta(|s_z|^2-(1-\lambda_z))\,d^4s_z=(1-\lambda_z)\,
\delta(|s'_z|^2-1)\,d^4s'_z$ (this is where the dimension $4$ of $\R^4$ enters), and $2J\langle s_u,O\,s_v\rangle=2J^\lambda_{uv}
\langle s'_u,O\,s'_v\rangle$; hence the factor $Z(\bs J^\lambda)/Z(J)$.
(v) Convergence: after $N$ applications of Lemma~\ref{lem:ipp} the identity holds with a remainder that is a finite sum over the
unfinished walks of length $N$. Each term still contains a linear factor of absolute value at most $1$ (all radii are at most $1$)
and $s_{y,b}$, and is bounded by $e^{4J|E_G|}J^N\prod_z1/(n_z-1)!$. On a graph with $|G|$ sites,
$\sum_z(n_z-1)\ge N+1-|G|$, so $\prod_z(n_z-1)!$ grows faster than any exponential of $N$, while there are at most $\Delta_G^N$ walks of
length $N$ ($\Delta_G$ the maximal degree). The remainder tends to $0$, and the finished terms are dominated in the same way.
\end{proof}

\begin{corollary}[matrix elements]\label{cor:matrix}
In slab $k$, for $x\ne y$, conditionally on $\cH$,
\[
\E\big[(Q^k_x)_{ij}\,((Q^k_y)^{-1})_{lk'}\,\big|\,\cH\big]=\sum_{\omega:x\to y}J^{|\omega|}\,(A_\omega)_{ik'}\,(B_\omega^{-1})_{lj}
\int F_\omega(\lambda)\,\frac{Z_k(\bs J^\lambda)}{Z_k(J)}\,d\lambda ,
\]
where $A_\omega$ (resp.\ $B_\omega$) is the ordered product of the horizontal links along $\omega$ in layer $k$ (resp.\ $k+1$).
\end{corollary}
\begin{proof}
Write $\iota(q)=\sum_aq_ae_a$ with $e_a\in M_2(\mathbb C)$. Then
\[
\sum_a(e_a)_{pq}\overline{(e_a)_{rs}}=2\delta_{pr}\delta_{qs},\qquad \sum_a(O_{A,B})_{ab}e_a=A\,e_bB^{-1},
\]
so $O_\omega=O_{A_\omega,B_\omega}$ and $\sum_{a,b}(O_\omega)_{ab}(e_a)_{ij}\overline{(e_b)_{k'l}}=2(A_\omega)_{ik'}
(B_\omega^{-1})_{lj}$. Insert \eqref{eq:walk}; the factor $2$ cancels the $\frac12$.
\end{proof}

\subsection{The Wilson loop as a sum over sequences of walks}
Let $\gamma$ have corners $(x,0)$, $(y,0)$, $(y,T)$, $(x,T)$ with $y=x+Re_1$; let $A$ be the product of the $R$ horizontal links from
$x$ to $y$ in layer $0$ and $B$ that from $y$ to $x$ in layer $T$. Then
$U_\gamma=A\,Q^0_y\cdots Q^{T-1}_y\,B\,(Q^{T-1}_x)^{-1}\cdots(Q^0_x)^{-1}$.

\begin{proposition}\label{prop:loop}
With walks $\omega_k:y\to x$ in $X$, one for each slab $k=0,\dots,T-1$,
\begin{equation}\label{eq:loop}
\langle W_\gamma\rangle_{L,\beta}=\frac12\sum_{\omega_0,\dots,\omega_{T-1}}\prod_{k<T}J^{|\omega_k|}\int\prod_{k<T}
F_{\omega_k}(\lambda^{(k)})\,d\lambda^{(k)}\ \E\Big[\prod_{k<T}\frac{Z_k(\bs J^{\lambda^{(k)}})}{Z_k(J)}\prod_{m=0}^{T}\Tr U_{C_m}\Big],
\end{equation}
where $\E$ is the expectation of the horizontal links under $\mu_{L,\beta}$ and $C_m$ is a closed loop of layer $m$: $C_0$ is the
bottom side ($x\to y$) followed by $\omega_0$ ($y\to x$); for $1\le m\le T-1$, $C_m$ is $\omega_{m-1}$ reversed ($x\to y$) followed by
$\omega_m$; $C_T$ is $\omega_{T-1}$ reversed followed by the top side ($y\to x$).
\end{proposition}
\begin{proof}
\emph{Contraction.} For matrices $P,N$ that do not depend on the vertical links of slab $k$, Corollary~\ref{cor:matrix} (walks from
$y$ to $x$) gives, conditionally on $\cH$ and on the other slabs,
\[
\E_k\big[\Tr\big(P\,Q^k_yN\,(Q^k_x)^{-1}\big)\big]=\sum_\omega w_\omega\,\Tr(PA'_\omega)\,\Tr(B'^{-1}_\omega N),
\]
with $w_\omega$ the weight of $\omega$ and $A'_\omega,B'_\omega$ its holonomies from $y$ to $x$ in layers $k$ and $k+1$: expand
$\Tr(PQ_yNQ_x^{-1})=\sum P_{ai}(Q_y)_{ij}N_{jl}(Q_x^{-1})_{la}$ and replace $(Q_y)_{ij}(Q_x^{-1})_{la}$ by
$(A'_\omega)_{ia}(B'^{-1}_\omega)_{lj}$. Apply this to slab $0$ with $P=A$, which produces
$\Tr U_{C_0}$ and leaves an expression of the same form with $P=B'^{-1}_{\omega_0}$; by induction over the slabs (independent given
$\cH$, Lemma~\ref{lem:slab}) one obtains $\prod_{m=0}^T\Tr U_{C_m}$.
\emph{Fubini.} For all $\lambda$ and all horizontal configurations, $0<Z_k(\bs J^\lambda)/Z_k(J)\le e^{12JL^3}$ (each of the $3L^3$
edges changes the exponent by at most $2J$), $|F_\omega|\le\prod_z1/(n_z-1)!$ and $|\Tr U_C|\le2$; the series
$\sum_\omega J^{|\omega|}\prod_z1/(n_z-1)!$ converges (proof of Proposition~\ref{prop:walk}). Sums, integrals and expectation may
therefore be exchanged.
\end{proof}

\subsection{Transfer matrix}\label{sec:transfer}
Let $\mathfrak h:=L^2(SU(2)^{E_X})$ with Haar measure. For $J\ge0$ let $t_J$ be the convolution operator
$(t_Jf)(v)=\int e^{J\ReTr(vw^{-1})}f(w)\,dw$ on $L^2(SU(2))$, and for $\bs J=(J_\ell)_{\ell\in E_X}$ let $T_0(\bs J):=\bigotimes_\ell
t_{J_\ell}$. Let $P_0$ be the average over gauge transformations $(u\cdot v)_{x,\mu}=u_xv_{x,\mu}u_{x+e_\mu}^{-1}$, $u\in SU(2)^X$
(the orthogonal projection onto gauge-invariant functions), and $M$ the multiplication by $\exp(J\sum_p\ReTr v_p)$, the sum running
over the $3L^3$ horizontal plaquettes of a layer. Then:
\begin{itemize}
\item $Z_k(\bs J)(v_k,v_{k+1})=\int\prod_xdu_x\prod_\ell e^{J_\ell\ReTr(v_{k,\ell}(u\cdot v_{k+1})_\ell^{-1})}$ is the kernel of
$T_0(\bs J)P_0$ (the vertical links of slab $k$ act as a gauge transformation of layer $k+1$);
\item $T_0(\bs J)$ commutes with gauge transformations (each $t_{J_\ell}$ is a convolution by a class function), hence with $P_0$;
\item $M$ and the functions $v\mapsto\Tr U_C(v)$ of closed loops $C$ are gauge invariant, so the corresponding multiplication
operators commute with $P_0$.
\end{itemize}
By Lemma~\ref{lem:pw} below, $t_J\ge0$ is trace class, and so is $T_0(\bs J)$. We put, for $\lambda\in[0,1]^{V(\omega)}$,
\[
\cT:=M^{1/2}P_0T_0(J)P_0M^{1/2}\ \ge0,\qquad \cT^{(\lambda)}:=M^{1/2}P_0T_0(\bs J^\lambda)P_0M^{1/2}\ \ge0 .
\]

\begin{proposition}[trace formula; {\cite[(B.10), p.~375]{BS83}}]\label{prop:trace}
For walks $\omega_0,\dots,\omega_{T-1}$ and parameters $\lambda^{(k)}$,
\begin{equation}\label{eq:trace}
\E\Big[\prod_{k<T}\frac{Z_k(\bs J^{\lambda^{(k)}})}{Z_k(J)}\prod_{m=0}^{T}\Tr U_{C_m}\Big]=\frac{\Tr\big(\prod_{m=0}^{L-1}\chi_m\cT_m
\big)}{\Tr\cT^L},
\end{equation}
where $\chi_m$ is the multiplication by $\Tr U_{C_m}$ for $m\le T$ and $\chi_m:=1$ for $m>T$, and $\cT_m:=\cT^{(\lambda^{(m)})}$ for
$m<T$, $\cT_m:=\cT$ for $m\ge T$ (the product is ordered by increasing $m$).
\end{proposition}
\begin{proof}
Each plaquette is either horizontal (in one layer) or vertical (in one slab), and the vertical links of slab $k$ are integrated in
$Z_k$; hence the marginal law of the horizontal links is $Z^{-1}\prod_mM(v_m)\prod_kZ_k(J)(v_k,v_{k+1})\prod dv$, with
$Z=\Tr\big((MT_0(J)P_0)^L\big)$. Multiplying by the ratios replaces the kernel of slab $k$ by that of $T_0(\bs J^{\lambda^{(k)}})P_0$,
and the left side of \eqref{eq:trace} becomes $Z^{-1}\Tr\big(\prod_m\chi_mMT_0(\bs J^{(m)})P_0\big)$. Write $M=M^{1/2}M^{1/2}$ and
$T_0P_0=P_0T_0P_0$; since $P_0$, $M^{1/2}$ and the $\chi_m$ commute, cyclicity of the trace gives $\Tr(\prod_m\chi_m\cT_m)$, and
likewise $Z=\Tr\cT^L$.
\end{proof}

\subsection{The bound of Borgs and Seiler}\label{sec:bsbound}
Borgs and Seiler note that $\cT^{(\lambda)}\le\cT$, since $\langle\varphi,\cT^{(\lambda)}\varphi\rangle$ is non-decreasing in each
coupling by reflection positivity \cite[(B.11), p.~376]{BS83}; Hölder's inequality then bounds the right side of \eqref{eq:trace} by
$2^{T+1}$ (Proposition~\ref{prop:holder} below with $c\le1$). Since $\int_0^1f_n=a_n:=1/(n+1)!$ \cite[(B.14), p.~376]{BS83}, \eqref{eq:loop}
gives
\begin{equation}\label{eq:bs}
|\langle W_\gamma\rangle|\le2^TG(x,y)^T,\qquad G(x,y):=\sum_{\omega:y\to x}J^{|\omega|}\prod_{z\in V(\omega)}a_{n_z(\omega)} .
\end{equation}
Read along the walk, $\prod_za_{n_z(\omega)}=\prod_{j=0}^{|\omega|}r(N_j)$ with $N_j:=\#\{i\le j:\omega_i=\omega_j\}$, $r(1)=a_1=\frac12$
and $r(N)=a_N/a_{N-1}=1/(N+1)$ for $N\ge2$. At each step after the first, the return to the previous site has $N\ge2$, so the six
continuations contribute at most $J(5\cdot\frac12+\frac13)=\frac{17}6J$ \cite[p.~376]{BS83}. Hence $G(x,y)\le\frac32\sum_{n\ge R}
J^n(\frac{17}6)^{n-1}$, and \eqref{eq:bs} gives the area law for $J<6/17$, i.e.\ $\beta<12/17$ \cite[p.~376]{BS83} \st{read}.

\begin{remark}[what Borgs and Seiler state]\label{rem:bs-statement}
Borgs and Seiler state the bound for the correlation of Polyakov loops at nonzero temperature: ``Irrespective of $J_M$ and
temperature we are sure to have confinement for $J_E<\tilde J_c$'', with $\tilde J_c=(d-\frac16)^{-1}$ for $SU(2)$
\cite[(II.55), p.~343]{BS83}, proved in their Appendix~B (pp.~373--377), uniformly in the number $L_0$ of time layers; they relate the Polyakov and
Wilson string tensions in \cite[(II.48)--(II.50), p.~341]{BS83} \st{read}. The version for rectangular Wilson loops at zero
temperature, with constants, is the adaptation written in \S\S\ref{sec:bs}--\ref{sec:bsbound}; in this note
``the threshold of the Borgs--Seiler method'' refers to the resulting range $\beta<12/17$. Page numbers of \cite{BS83} refer to the
journal version; the preprint MPI-PAE/PTh 16/83 has the same text and equation numbers in all the passages cited.
\end{remark}

\section{Improvement (I): a strict operator inequality}\label{sec:I}

\begin{lemma}[Peter--Weyl spectrum]\label{lem:pw}
On the span of the matrix coefficients of the spin-$j$ representation, $j\in\{0,\frac12,1,\dots\}$, the operator $t_J$ acts as
multiplication by
\[
m_j(J)=\frac{I_{2j+1}(2J)}{J}=\sum_{k\ge0}\frac{J^{2k+2j}}{k!\,(k+2j+1)!}\quad(J>0),\qquad m_0(0)=1,\ m_j(0)=0\ (j\ge\tfrac12).
\]
In particular $t_J\ge0$ and $t_J$ is trace class. For $0<J\le1$ and $0\le t\le1$, $m_j(tJ)/m_j(J)\le m_0(tJ)/m_0(J)$ for every
$j$.
\end{lemma}
\begin{proof}
Convolution by the class function $\varphi(g)=e^{J\Tr g}$ acts on the spin-$j$ component by $c_j/(2j+1)$, $c_j:=\int\varphi\,\chi_j$.
With $\Tr g=2\cos\vartheta$, $\chi_j=\sin((2j+1)\vartheta)/\sin\vartheta$ and Weyl's integration formula,
$c_j=\frac2\pi\int_0^\pi e^{2J\cos\vartheta}\sin((2j+1)\vartheta)\sin\vartheta\,d\vartheta=I_{2j}(2J)-I_{2j+2}(2J)=(2j+1)I_{2j+1}(2J)/J$,
and the series is that of $I_{2j+1}$. The eigenvalues are positive and $\sum_j(2j+1)^2m_j(J)<\infty$.
For $j\ge\frac12$, termwise comparison gives $m_j(tJ)\le t^{2j}m_j(J)\le t\,m_j(J)$. For $j=0$, $m_0(tJ)/m_0(J)=\sum_kw_kt^{2k}$ with
probability weights $w_k\propto J^{2k}/(k!(k+1)!)$; by Jensen's inequality this is at least $t^{2\langle k\rangle}$, and
$2\langle k\rangle\le1$ because $\sum_{k\ge1}(2k-1)J^{2k}/(k!(k+1)!)\le\sum_{k\ge1}(2k-1)/(k!(k+1)!)<0.79$ for $J\le1$. Hence
$m_0(tJ)/m_0(J)\ge t$. (For $t=0$ the left side vanishes for $j\ge\frac12$ and equals $1/m_0(J)$ for $j=0$.)
\end{proof}

For $s\in[0,1]$ put
\begin{equation}\label{eq:rho}
M_J(s):=\sum_{k\ge0}\frac{(J^2s)^k}{k!\,(k+1)!}=m_0(J\sqrt s),\qquad \rho_J(s):=\frac{M_J(s)}{M_J(1)}\in(0,1].
\end{equation}

\begin{proposition}[strict operator inequality]\label{prop:strict}
Let $0<J\le1$, $\omega$ a walk in $X$ and $\lambda\in[0,1]^{V(\omega)}$. Then $0\le\cT^{(\lambda)}\le c(\lambda)\,\cT$ with
\[
c(\lambda):=\prod_{\{u,v\}\in E_X}\rho_J\big((1-\lambda_u)(1-\lambda_v)\big)\qquad(\lambda_u:=0\text{ for }u\notin V(\omega)).
\]
The constant is optimal link by link.
\end{proposition}
\begin{proof}
$T_0(\bs J^\lambda)$ and $T_0(J)$ are diagonal in the product Peter--Weyl basis, with eigenvalues $\prod_\ell m_{j_\ell}(J^\lambda_\ell)$
and $\prod_\ell m_{j_\ell}(J)$. Since $J^\lambda_\ell=J\sqrt{s_\ell}$ with $s_\ell=(1-\lambda_u)(1-\lambda_v)$, Lemma~\ref{lem:pw}
gives $T_0(\bs J^\lambda)\le c(\lambda)T_0(J)$; conjugation by $M^{1/2}P_0$ preserves the order. The constant function (all
$j_\ell=0$) is gauge invariant and gives equality for $T_0$.
\end{proof}
\noindent Borgs and Seiler use $\cT^{(\lambda)}\le\cT$, i.e.\ $c\le1$. Here $c(\lambda)<1$ as soon as $\lambda\not\equiv0$. This is
improvement~(I). (For $J=0$ the operators $\cT^{(\lambda)}$ and $\cT$ coincide, $c\equiv1$ and $\langle W_\gamma\rangle=0$: the case is
trivial, and it is excluded from the statements above, as in \S\ref{sec:proof}.)

\begin{proposition}[Hölder bound]\label{prop:holder}
In the setting of Proposition~\ref{prop:trace}, with $c_k:=c(\lambda^{(k)})$,
\[
\Big|\E\Big[\prod_{k<T}\frac{Z_k(\bs J^{\lambda^{(k)}})}{Z_k(J)}\prod_{m=0}^T\Tr U_{C_m}\Big]\Big|\le2^{T+1}\prod_{k<T}c_k .
\]
\end{proposition}
\begin{proof}
If $0\le A\le B$ are bounded operators, then $\|A^{1/2}f\|\le\|B^{1/2}f\|$, so there is a contraction $C$ with $A^{1/2}=CB^{1/2}$
\cite{Dou66} and $A=B^{1/2}C^*CB^{1/2}$. With $A=\cT_k$ and $B=c_k\cT$ this gives $\cT_k=\cT^{1/2}Y_k\cT^{1/2}$ with $0\le Y_k$,
$\|Y_k\|\le c_k$ (and $Y_m=I$ for $m\ge T$). By cyclicity, $\Tr(\prod_m\chi_m\cT_m)=\Tr(\prod_mX_mY_m)$ with
$X_m:=\cT^{1/2}\chi_m\cT^{1/2}$. Hölder's inequality for Schatten norms, with $L$ factors of exponent $L$ and the $Y_m$ in operator
norm, gives $|\Tr(\cdots)|\le\prod_m\|X_m\|_L\|Y_m\|$, and $\|X_m\|_L\le\|\chi_m\|_\infty\|\cT^{1/2}\|_{2L}^2=\|\chi_m\|_\infty
(\Tr\cT^L)^{1/L}$. Conclude with $\|\chi_m\|_\infty\le2$ for $m\le T$, $=1$ otherwise, and \eqref{eq:trace}.
\end{proof}

\begin{lemma}[Cauchy--Schwarz]\label{lem:cs}
For $a,b\in[0,1]$: $M_J(ab)^2\le M_J(a^2)M_J(b^2)$ and $M_J(a)\le M_J(a^2)^{1/2}M_J(1)^{1/2}$.
\end{lemma}
\begin{proof}
Cauchy--Schwarz for $\sum_kc_k(a^k)(b^k)$ with $c_k=J^{2k}/(k!(k+1)!)\ge0$; the second inequality is the case $b=1$.
\end{proof}

\begin{corollary}[site-by-site factorisation]\label{cor:site}
Let $L\ge3$, $\omega$ a walk in $X$ and $\lambda\in[0,1]^{V(\omega)}$. Then
\[
c(\lambda)\le\prod_{z\in V(\omega)}h_J(\lambda_z),\qquad h_J(\lambda):=\rho_J\big((1-\lambda)^2\big)^3 .
\]
\end{corollary}
\begin{proof}
Put $a_z:=1-\lambda_z$. By Lemma~\ref{lem:cs}, an edge $\{u,v\}\subset V(\omega)$ contributes $\rho(a_ua_v)\le\rho(a_u^2)^{1/2}
\rho(a_v^2)^{1/2}$, an edge $\{u,w\}$ with $w\notin V(\omega)$ contributes $\rho(a_u)\le\rho(a_u^2)^{1/2}$, and the other edges
contribute $1$. Every site of $X$ has exactly six incident edges, so every $z\in V(\omega)$ receives exactly six factors
$\rho(a_z^2)^{1/2}$.
\end{proof}
\noindent The exponent $3$ uses that every site has exactly six edges, which holds on the torus with $L\ge3$. A site with $k<6$ edges
(for instance on the boundary of a box) would receive $\rho(a_z^2)^{k/2}\ge\rho(a_z^2)^3$; this is why the theorem is stated on tori
only (\S\ref{sec:scope}).

\paragraph{New weights.} Let
\begin{equation}\label{eq:At}
\At_n(J):=\int_0^1f_n(\lambda)\,h_J(\lambda)\,d\lambda\ \le\ a_n=\frac1{(n+1)!},\qquad \rt(1):=\At_1,\quad \rt(N):=\frac{\At_N}{\At_{N-1}}\
(N\ge2).
\end{equation}
Combining \eqref{eq:loop}, Proposition~\ref{prop:holder} and Corollary~\ref{cor:site} (all integrands are non-negative after the
bound), and integrating site by site,
\begin{equation}\label{eq:Gt}
|\langle W_\gamma\rangle_{L,\beta}|\le2^T\,\Gt(x,y)^T,\qquad \Gt(x,y):=\sum_{\omega:y\to x}J^{|\omega|}\prod_{z\in V(\omega)}
\At_{n_z(\omega)}=\sum_{\omega:y\to x}J^{|\omega|}\prod_{j=0}^{|\omega|}\rt(N_j),
\end{equation}
with $N_j=\#\{i\le j:\omega_i=\omega_j\}$ as in \S\ref{sec:bsbound}: the weight is read step by step.

\begin{lemma}[revisits]\label{lem:revisit}
$\rt(N)\le1/(N+1)$ for $N\ge2$.
\end{lemma}
\begin{proof}
$\rt(N)=\E_\nu[\lambda]/(N-1)$ with $\nu\propto\lambda^{N-2}(1-\lambda)h_J(\lambda)\,d\lambda$. The function $h_J$ is non-increasing
($\rho_J$ is non-decreasing and $(1-\lambda)^2$ decreasing). By Chebyshev's integral inequality for the positive measure
$w=\lambda^{N-2}(1-\lambda)\,d\lambda$, the increasing function $\lambda$ and the non-increasing function $h_J$,
$\E_\nu[\lambda]=\int w\lambda h_J/\int wh_J\le\int w\lambda/\int w=(N-1)/(N+1)$.
\end{proof}

\begin{remark}
With the count of \S\ref{sec:bsbound} (memory $1$), (I) alone gives the range $\beta\lesssim0.7824$ \st{numerical}. At $J=J_*$ one
has $\At_1=0.443739\ldots$ instead of $a_1=\frac12$.
\end{remark}

\section{Improvement (II): walks with memory, and proof of Theorem~\ref{thm:main}}\label{sec:II}

\subsection{Monotonicity in the coupling}
\begin{lemma}\label{lem:monotone}
(M1) For every $\lambda$, $J\mapsto h_J(\lambda)$ is non-increasing on $(0,\infty)$. (M2) For $0<J\le0.44$, $J\mapsto J\,h_J(\lambda)$
is non-decreasing. Consequently, for $0<J\le J'\le0.44$ and every $N\ge1$, $J\,\rt(N;J)\le J'\,\rt(N;J')$.
\end{lemma}
\begin{proof}
Let $F(x):=\sum_kx^k/(k!(k+1)!)$ and $\varphi(x):=xF'(x)/F(x)$, the mean of $k$ under the weights $x^k/(k!(k+1)!)$, which is
non-decreasing in $x$ (its logarithmic derivative is a variance). Then $\partial_J\log\rho_J(s)=\frac2J(\varphi(J^2s)-\varphi(J^2))\in
[-\frac2J\varphi(J^2),0]$, so $\partial_J\log h_J\in[-\frac6J\varphi(J^2),0]$, which is (M1), and $\partial_J\log(Jh_J)\ge\frac1J
(1-6\varphi(J^2))$. For $J\le0.44$, $J^2\le0.2$ and, since $F\ge1$,
$\varphi(J^2)\le0.2\,F'(0.2)=\sum_{k\ge1}k\,(0.2)^k/(k!(k+1)!)=0.10683557\ldots<\frac16$: the first $40$ terms are summed in exact
rational arithmetic, and the tail is at most $t_{41}/(1-0.2/(41\cdot43))<2.5\cdot10^{-128}$, where $t_k:=k(0.2)^k/(k!(k+1)!)$ satisfies
$t_{k+1}/t_k=0.2/(k(k+2))$. This proves (M2). Hence $J\At_N(J)=\int f_N\cdot(Jh_J)$ is non-decreasing and $\At_{N-1}(J)$ is
non-increasing, so $J\rt(N;J)$ is non-decreasing ($N=1$: $J\At_1$).
\end{proof}

\subsection{Counting with a finite memory}
Let $D:=\{\pm e_1,\pm e_2,\pm e_3\}$ and write a walk through its steps $d_j:=\omega_j-\omega_{j-1}\in D$. Fix $m\ge1$. For $j\ge m+1$
the \emph{state} before step $j$ is $s=(s_1,\dots,s_m):=(d_{j-1},\dots,d_{j-m})\in D^m$, and the \emph{window count} is
\[
N_j^{(m)}:=1+\#\{1\le i\le m:\ \omega_{j-1-i}=\omega_j\}=N(s,d_j),
\]
where $N(s,d):=1+\#\{1\le i\le m:\ -(s_1+\dots+s_i)=d\}$ (indeed $\omega_{j-1-i}-\omega_{j-1}=-(s_1+\dots+s_i)$),
so that $N_j^{(m)}\le N_j$. Let $\rb:\N_{\ge1}\to(0,\infty)$ be non-increasing with $\rb(N)\ge\sup_{N'\ge N}\rt(N')$, and let $T_m$ be
the operator on functions on $D^m$
\[
(T_mf)(s):=\sum_{d\in D}\rb\big(N(s,d)\big)\,f(d,s_1,\dots,s_{m-1}).
\]

\begin{lemma}[window]\label{lem:window}
If $H:D^m\to(0,\infty)$ satisfies $T_mH\le\rho'H$ pointwise, then for every $n\ge m$ and every starting point,
\[
\sum_{|\omega|=n}\ \prod_{j=1}^{n}\rt(N_j)\ \le\ (6\,\rb(1))^m\,\frac{H_{\max}}{H_{\min}}\,\rho'^{\,n-m},
\]
the sum running over all walks of length $n$ in $\Z^3$ from that point.
\end{lemma}
\begin{proof}
For $j\le m$ use $\rt(N_j)\le\rb(1)$; the $6^m$ choices of $(d_1,\dots,d_m)$ give $(6\rb(1))^m$ and run over all states. For $j>m$,
$\rt(N_j)\le\sup_{N'\ge N_j^{(m)}}\rt(N')\le\rb(N(s,d_j))$. Summing the last $n-m$ steps gives $(T_m^{n-m}\mathbf 1)(s)\le
(T_m^{n-m}H)(s)/H_{\min}\le\rho'^{\,n-m}H_{\max}/H_{\min}$, since $T_m$ has non-negative coefficients.
\end{proof}

The hyperoctahedral group $\mathfrak G$ (the $48$ signed permutations of the axes) acts on $D$ and diagonally on $D^m$; since
$N(gs,gd)=N(s,d)$ and $(d,s_1,\dots,s_{m-1})$ is equivariant, $T_m$ commutes with $\mathfrak G$. For a $\mathfrak G$-invariant $H$
the inequality $T_mH\le\rho'H$ holds everywhere as soon as it holds at one state of each orbit.

\begin{lemma}[torus]\label{lem:torus}
Let $L\ge3$. The walks of length $n$ in $X=(\Z/L\Z)^3$ from a site are in bijection with $D^n$, i.e.\ with the walks of length $n$ in
$\Z^3$ from a lift of that site, and the visit counts on the torus, global and in windows, are at least those of the lift. Hence the
bound of Lemma~\ref{lem:window} holds for walks on the torus. A walk from $y$ to $x=y-Re_1$ in $X$ has length at least $R$ when
$R\le L/2$.
\end{lemma}
\begin{proof}
Coincidences in $\Z^3$ project to coincidences in $X$. For $j>m$, $\rt(N_j^{X})\le\rb(N_j^{X,(m)})\le\rb(N_j^{\Z^3,(m)})$ because $\rb$
is non-increasing. The torus distance of $x$ and $y$ is $\min(R,L-R)=R$.
\end{proof}

\subsection{The certificate}
\begin{proposition}[certificate, $m=11$]\label{prop:cert}
Let $J_*=40674/100000$ and let $p_1,\dots,p_{13}$ be the integers of the weights file of \S\ref{sec:repro}
($p_1=487\,896\,913\,543$, \dots, $p_{13}=78\,536\,544\,842$); put $\rb_*(N):=p_{\min(N,13)}/2^{40}$. Then:
\begin{enumerate}
\item $\rb_*$ is non-increasing and $\rb_*(N)\ge\sup_{N'\ge N}\rt(N';J_*)$ for every $N\ge1$ \st{certified};
\item there is a $\mathfrak G$-invariant $H:D^{11}\to\N$ with $H_{\min}=663\,930\,282\,413$ and $H_{\max}=2^{40}$ such that, for every
$s\in D^{11}$,
\begin{equation}\label{eq:cert}
2^{30}\sum_{d\in D}p_{\min(N(s,d),13)}\,H(d,s_1,\dots,s_{10})\ \le\ 2639854083\cdot2^{40}\cdot H(s)
\end{equation}
\st{certified}, i.e.\ $T_{11}H\le\rho'_*H$ for the weights $\rb_*$;
\item $\theta_*=J_*\rho'_*=53686712485971/53687091200000<1$ \st{certified}.
\end{enumerate}
\end{proposition}
\noindent The computations and their verification are described in \S\ref{sec:cap}.

\subsection{Proof of Theorem~\ref{thm:main}}\label{sec:proof}
Let $\beta\in[0,\beta_*]$, so $J=\beta/2\le J_*\le0.44$; for $J=0$ the loop expectation vanishes, so let $J>0$. By \eqref{eq:Gt} it
suffices to bound $\Gt(x,y)=\sum_nJ^nS_n$, $S_n:=\sum_{\omega:y\to x,\,|\omega|=n}\prod_{j=0}^n\rt(N_j;J)$. By
Lemma~\ref{lem:monotone}, $J\rt(N;J)\le J_*\rt(N;J_*)$ for all $N$, so for every walk
$J^n\prod_{j=0}^n\rt(N_j;J)\le\At_1(J)\prod_{j=1}^nJ_*\rt(N_j;J_*)$; only the exact weights are compared, never the rounded ones. By
Lemmas~\ref{lem:window} and~\ref{lem:torus} with $\rt(\cdot;J_*)$, $\rb_*$, $H$ and $\rho'_*$ (Proposition~\ref{prop:cert}), for
$n\ge m=11$,
\[
J^nS_n\le\At_1(J)\,(6J_*\rb_*(1))^m\,\frac{H_{\max}}{H_{\min}}\,(J_*\rho'_*)^{n-m}=\At_1(J)\,K_0\,\theta_*^{\,n},
\]
\begin{equation}\label{eq:K0}
K_0:=\Big(\frac{6\,\rb_*(1)}{\rho'_*}\Big)^{11}\frac{H_{\max}}{H_{\min}}=3.978058961\ldots .
\end{equation}
For $n<m$, bounding each step by $6J_*\rb_*(1)$ gives $J^nS_n\le\At_1(J)(6\rb_*(1)/\rho'_*)^n\theta_*^n\le\At_1(J)K_0\theta_*^n$,
because $6\rb_*(1)/\rho'_*=1.0829\ldots>1$ and $H_{\max}/H_{\min}\ge1$. Walks from $y$ to $x$ have length $n\ge R$
(Lemma~\ref{lem:torus}) and $\At_1(J)\le a_1=\frac12$, so $\Gt(x,y)\le\frac12K_0\,\theta_*^R/(1-\theta_*)$, and \eqref{eq:Gt} gives
$|\langle W_\gamma\rangle|\le(K_0/(1-\theta_*))^T\theta_*^{RT}$. \qed

\begin{remark}[the two improvements separately, and the role of $m$]\label{rem:contrib}
With the weights $r(N)$ of Borgs and Seiler, the same operators give the spectral radii $\rho_1=17/6$ (the bound of
\S\ref{sec:bsbound}), $\rho_3=2.7869954984$, $\rho_5=2.7653748962$, $\rho_7=2.7528722513$, $\rho_9=2.7447464071$,
$\rho_{11}=2.7390488353$; (II) alone therefore gives $\beta<2/\rho_{11}=0.73018$ \st{numerical}. With the weights of (I),
Table~\ref{tab:memory} lists the thresholds; the increments decrease slowly with $m$ \st{numerical}.
\end{remark}

\begin{table}[ht]
\centering\small
\begin{tabular}{@{}rrllll@{}}
\toprule
$m$ & orbits & $J_*$ & $\beta_*$ & $\rho'_*$ & status\\
\midrule
$1$ & $1$ & & $0.7824$ & & numerical\\
$3$ & $11$ & & $0.7972$ & & numerical\\
$5$ & $236$ & & $0.8044$ & & numerical\\
$7$ & $6\,896$ & $40432/10^5$ & $0.80864$ & $2655621200/2^{30}$ & certified\\
$9$ & $226\,496$ & $40574/10^5$ & $0.81148$ & $2646370548/2^{30}$ & certified\\
$11$ & $7\,821\,056$ & $40674/10^5$ & $0.81348$ & $2639854083/2^{30}$ & certified (Theorem~\ref{thm:main})\\
$13$ & $276\,294\,656$ & $40749/10^5$ & $0.81498$ & $2635005744/2^{30}$ & certified (Remark~\ref{rem:m13})\\
\bottomrule
\end{tabular}
\caption{Thresholds of the combined method as a function of the memory $m$. For each certified row, $J_*+10^{-5}$ fails the test
$J\rho'<1$ of the dichotomy (grid $10^{-5}$ in $J$). The numbers of orbits are the Burnside counts $(6^m+3\cdot4^m+15\cdot2^m)/48$.}
\label{tab:memory}
\end{table}

\subsection{A tilted bound (new in version~1.1)}\label{sec:tilted}
The proof of Theorem~\ref{thm:main} bounds $\Gt(x,y)$ by a sum over \emph{all} walks of length $n\ge R$ from $y$, whatever their end
point. Near the threshold this costs the factor $1/(1-\theta_*)$, and the rate obtained, $-\log\theta_*$, is of order $1-\theta_*$.
Weighting each step by $e^{\eta\langle\hat e,d\rangle}$, where $\hat e$ is the direction from $y$ to $x$ (as for the tilted walk in the
proof of \cite[Thm.~7.1]{D7}), keeps the threshold and gives a rate of order $\sqrt{1-\theta_*}$ \st{numerical}, as Table~\ref{tab:tilted}
shows; the parity of walks in $\Z^3$ halves the constant on even tori. The weights, the window counts and Lemmas~\ref{lem:monotone},
\ref{lem:window} and~\ref{lem:torus} are those of \S\ref{sec:II}; only the transfer operator and the certificate change. We write $\eta$
for the tilt and $\sigma_\eta$ for the rate, to avoid confusion with the parameters $\lambda_z$ of the walk expansion and with the string
tension $\sigma(\beta)$.

For $\hat e\in D$ and $\eta\ge0$, let $T_m^{(\eta,\hat e)}$ act on functions on $D^m$ by
\[
(T_m^{(\eta,\hat e)}f)(s):=\sum_{d\in D}\rb\big(N(s,d)\big)\,e^{\eta\langle\hat e,d\rangle}\,f(d,s_1,\dots,s_{m-1}).
\]
Lemma~\ref{lem:window} holds for $T_m^{(\eta,\hat e)}$ with the same proof: if $T_m^{(\eta,\hat e)}H\le\rho'H$, the sum over the walks of
length $n\ge m$ from a point of $\prod_{j=1}^{n}\rt(N_j)\,e^{\eta\langle\hat e,d_j\rangle}$ is at most
$\big(\rb(1)(e^{\eta}+e^{-\eta}+4)\big)^m(H_{\max}/H_{\min})\,\rho'^{\,n-m}$, because $\sum_{d\in D}e^{\eta\langle\hat e,d\rangle}=e^\eta+e^{-\eta}+4$;
and Lemma~\ref{lem:torus} transfers this bound to tori. For $g\in\mathfrak G$ one has $N(gs,gd)=N(s,d)$ ($N$ only involves the equalities
$-(s_1+\dots+s_i)=d$, which the linear map $g$ preserves) and $\langle g\hat e,gd\rangle=\langle\hat e,d\rangle$; so $T_m^{(\eta,g\hat e)}$ is
conjugate to $T_m^{(\eta,\hat e)}$ by $s\mapsto gs$, and if $T_m^{(\eta,e_1)}H\le\rho'H$ then $T_m^{(\eta,ge_1)}(H\circ g^{-1})\le
\rho'\,H\circ g^{-1}$, with the same ratio $H_{\max}/H_{\min}$.

\begin{theorem}[tilted bound]\label{thm:tilted}
Let $m\ge1$ and $0<J_i\le0.44$; let $\rb$ be non-increasing with $\rb(N)\ge\sup_{N'\ge N}\rt(N';J_i)$ for all $N\ge1$; let $\eta\ge0$; and
let $H:D^m\to(0,\infty)$ and $\rho'>0$ satisfy $T_m^{(\eta,e_1)}H\le\rho'H$. Put
\[
\theta:=J_i\rho',\qquad \sigma_\eta:=\eta-\log\theta,\qquad q:=\frac{\rb(1)\,(e^{\eta}+e^{-\eta}+4)}{\rho'},
\]
\[
K_0^{(\eta)}:=\max(1,q)^{m}\,\frac{H_{\max}}{H_{\min}},\qquad C^{(\eta)}:=\frac{K_0^{(\eta)}}{1-\theta^{2}},\qquad
C_{\rm odd}^{(\eta)}:=\frac{K_0^{(\eta)}}{1-\theta},
\]
and assume $\theta<1$. Let $0\le\beta\le2J_i$, and let $\gamma$ be a rectangular loop with sides $R$ along $e_\mu$ and $T$ along $e_\nu$.
\begin{enumerate}
\item On the torus $(\Z/L\Z)^4$ with $L$ even and $L\ge\max(4,2R,T+1)$:
\[
|\langle W_\gamma\rangle_{L,\beta}|\ \le\ \big[C^{(\eta)}\big(1+e^{-\sigma_\eta(L-2R)}\big)\big]^{T}e^{-\sigma_\eta RT}\ \le\
\big(2C^{(\eta)}\big)^{T}e^{-\sigma_\eta RT}.
\]
\item On the torus with $L$ odd and $L\ge\max(3,2R+1,T+1)$: the same bound with $C_{\rm odd}^{(\eta)}$ in place of $C^{(\eta)}$.
\item Every weak limit point of $(\mu_{L,\beta})$ along even $L\to\infty$ satisfies
$|\langle W_\gamma\rangle|\le\big(C^{(\eta)}\big)^{T}e^{-\sigma_\eta RT}$.
\end{enumerate}
\end{theorem}
\noindent Status: \st{proved}, given $H$ and $\rho'$; the instances of Table~\ref{tab:tilted} are \st{certified}
(Proposition~\ref{prop:tiltcert}). Not formalised in Lean.

\begin{proof}
Let $J=\beta/2\in(0,J_i]$ (for $J=0$ the expectation vanishes), and let $\hat e\in D$ be the direction from $y$ to $x$ in \eqref{eq:Gt},
so that $x=y+R\hat e$. As in \S\ref{sec:proof}, Lemma~\ref{lem:monotone} applied to the exact weights gives, for every walk of length $n$,
$J^{n}\prod_{j=0}^{n}\rt(N_j;J)\le\At_1(J)\prod_{j=1}^{n}J_i\,\rt(N_j;J_i)$.

\emph{Walks in $\Z^3$.} A walk from $y$ to $y+R\hat e$ satisfies $\prod_{j=1}^{n}e^{\eta\langle\hat e,d_j\rangle}=e^{\eta R}$, $n\ge R$ and
$n\equiv R\pmod2$, since $\Z^3$ is bipartite. Multiplying each term by $e^{-\eta R}\prod_je^{\eta\langle\hat e,d_j\rangle}=1$ and dropping
the constraint on the end point (all terms are non-negative),
\begin{equation}\label{eq:tilt1}
\sum_{\omega:y\to y+R\hat e}J^{|\omega|}\prod_{j=0}^{|\omega|}\rt(N_j;J)\ \le\ \At_1(J)\,e^{-\eta R}\sum_{\substack{n\ge R\\ n\equiv R\,(2)}}
J_i^{\,n}\sum_{|\omega|=n}\ \prod_{j=1}^{n}\rt(N_j;J_i)\,e^{\eta\langle\hat e,d_j\rangle},
\end{equation}
the inner sum running over all walks of length $n$ from $y$. Write $\hat e=ge_1$ with $g\in\mathfrak G$. By the tilted form of
Lemma~\ref{lem:window} with the vector $H\circ g^{-1}$, $J_i^{\,n}\sum_{|\omega|=n}\cdots\le J_i^{\,n}(q\rho')^m(H_{\max}/H_{\min})\rho'^{\,n-m}=
q^m(H_{\max}/H_{\min})\,\theta^n\le K_0^{(\eta)}\theta^n$ for $n\ge m$. For $n<m$, bounding every step by $\rb(1)(e^\eta+e^{-\eta}+4)=q\rho'$
gives $q^n\theta^n\le\max(1,q)^m\theta^n\le K_0^{(\eta)}\theta^n$. Summing the geometric series over $n\ge R$, $n\equiv R$, and using
$\At_1\le\frac12$, the right side of \eqref{eq:tilt1} is at most $\frac12K_0^{(\eta)}e^{-\eta R}\theta^R/(1-\theta^2)=
\frac12C^{(\eta)}e^{-\sigma_\eta R}$.

\emph{Tori.} By Lemma~\ref{lem:torus}, a walk from $y$ to $x$ in $X=(\Z/L\Z)^3$ lifts to a unique walk in $\Z^3$ from a fixed lift of
$y$, ending at a lift of $x$, i.e.\ with displacement $R\hat e+Lk$, $k\in\Z^3$, and the window bound applies with the visit counts of the
lift (the counts on the torus are larger and $\rb$ is non-increasing). If $\langle\hat e,k\rangle\ge0$, the displacement along $\hat e$ is
at least $R$, so $e^{-\eta R}\prod_je^{\eta\langle\hat e,d_j\rangle}\ge1$ and these walks are bounded as in \eqref{eq:tilt1}. If
$\langle\hat e,k\rangle\le-1$, the displacement along $\hat e$ is $R+L\langle\hat e,k\rangle\le R-L\le-R$, so the length is at least $L-R$
and $e^{-\eta(L-R)}\prod_je^{\eta\langle-\hat e,d_j\rangle}=e^{-\eta(L-R)-\eta(R+L\langle\hat e,k\rangle)}\ge1$; inserting this factor and using
$T_m^{(\eta,-\hat e)}$, the same argument bounds these walks by $\frac12C^{(\eta)}e^{-\sigma_\eta(L-R)}$. For even $L$ the coordinates of
$Lk$ are even, so $n\equiv R\pmod2$ for every lift and $L-R\equiv R$. Therefore $\Gt(x,y)\le\frac12C^{(\eta)}e^{-\sigma_\eta R}\big(1+
e^{-\sigma_\eta(L-2R)}\big)$, and \eqref{eq:Gt} gives item~1. For odd $L$ the parity is not used, $\sum_{n\ge N}\theta^n=\theta^N/(1-\theta)$,
and the same argument gives item~2. Item~3 follows from item~1, since $e^{-\sigma_\eta(L-2R)}\to0$ and $W_\gamma$ is a bounded continuous
function of finitely many links.
\end{proof}

\begin{proposition}[tilted certificates, $m=11$]\label{prop:tiltcert}
For each row of Table~\ref{tab:tilted}, with $J_i=\beta_i/2$, the weights $\rb(N)=p_{\min(N,13)}/2^{40}$ of the weights file at $J_i$
(for $J_*$ that of Proposition~\ref{prop:cert}, for the other rows those of Table~\ref{tab:tension}), $e^\eta=a/b$ and the stated
integer $P$, there is $H:D^{11}\to\N$ with $H_{\max}=2^{40}$ and the stated $H_{\min}$ such that, for every $s\in D^{11}$,
\begin{equation}\label{eq:tiltcert}
\sum_{d\in D}c_d\,p_{\min(N(s,d),13)}\,H(d,s_1,\dots,s_{10})\ \le\ P\cdot2^{10}\,ab\,H(s),
\end{equation}
where $c_{+e_1}=a^2$, $c_{-e_1}=b^2$ and $c_d=ab$ for the four other steps $d$; that is, $T_{11}^{(\eta,e_1)}H\le(P/2^{30})H$ \st{certified}; moreover $\theta=J_iP/2^{30}<1$ \st{certified}. Consequently
Theorem~\ref{thm:tilted} holds with $\rho'=P/2^{30}$ and the constants of Table~\ref{tab:tilted}.
\end{proposition}

\begin{table}[ht]
\centering\footnotesize\setlength{\tabcolsep}{3.5pt}
\begin{tabular}{@{}lllllllll@{}}
\toprule
$\beta\le$ & $e^\eta$ & $P$ & $H_{\min}$ & $1-\theta$ & $\sigma_\eta\ge$ & $C^{(\eta)}\le$ & $R_{\min}^{\rm lim}\le$ & $R_{\min}^{\rm tor}\le$\\
\midrule
$0.81348$ & $201/200$ & $2\,639\,868\,799$ & $663\,255\,845\,757$ & $1.4795\cdot10^{-6}$ & $0.004989$ & $1.3457\cdot10^{6}$ & $2\,829$ & $3\,107$\\
$0.81348$ & $251/250$ & $2\,639\,863\,512$ & $663\,389\,751\,364$ & $3.4823\cdot10^{-6}$ & $0.003995$ & $5.7164\cdot10^{5}$ & $3\,318$ & $3\,665$\\
$0.81$ & $10/9$ & $2\,648\,763\,350$ & $647\,937\,564\,146$ & $9.2449\cdot10^{-4}$ & $0.1062$ & $2\,192.6$ & $72.4$ & $85.5$\\
$0.80$ & $5/4$ & $2\,678\,382\,511$ & $632\,089\,802\,916$ & $2.2247\cdot10^{-3}$ & $0.2253$ & $915.6$ & $30.3$ & $36.5$\\
$0.70588$ & $7/4$ & $2\,898\,427\,727$ & $568\,404\,757\,891$ & $4.7283\cdot10^{-2}$ & $0.6080$ & $43.8$ & $6.3$ & $8.5$\\
\bottomrule
\end{tabular}
\caption{Tilted certificates with memory $11$ (full state space $D^{11}$, no symmetry reduction) and the constants of
Theorem~\ref{thm:tilted}. The exact $\theta$ is the rational $J_iP/2^{30}$; $1-\theta$ and $\sigma_\eta$ are rounded down; $C^{(\eta)}$,
$R_{\min}^{\rm lim}:=\log C^{(\eta)}/\sigma_\eta$ (weak limits) and $R_{\min}^{\rm tor}:=\log(2C_{\rm odd}^{(\eta)})/\sigma_\eta$ (every torus,
even or odd) are rounded up; beyond these radii the bounds are below $1$. For comparison, Theorem~\ref{thm:main} and Table~\ref{tab:tension}
give $\sigma\ge7.05\cdot10^{-6}$, $0.0034$, $0.0133$, $0.1159$ and $R_{\min}\approx1.9\cdot10^{6}$, $2.1\cdot10^{3}$, $430$, $31$ (every torus) for
$\beta\le0.81348$, $0.81$, $0.80$, $0.70588$.}
\label{tab:tilted}
\end{table}

\begin{corollary}[string tension, version~1.1]\label{cor:tilted}
With $\sigma(\beta)$ as in Corollary~\ref{cor:tension}: $\sigma(\beta)\ge0.004989$ for $0<\beta\le\beta_*$, $\sigma(\beta)\ge0.1062$ for
$0<\beta\le0.81$, $\sigma(\beta)\ge0.2253$ for $0<\beta\le0.80$, and $\sigma(\beta)\ge0.6080$ for $0<\beta\le0.70588$.
\end{corollary}
\begin{proof}
Items~1 and~2 of Theorem~\ref{thm:tilted} cover every $L\ge\max(3,2R,T+1)$ (for odd $L$, $L\ge2R$ is the same as $L\ge2R+1$), with the
constant $C':=2C_{\rm odd}^{(\eta)}\ge2C^{(\eta)}$, which depends neither on $L$ nor on $T$. Hence
$-\frac1{RT}\log|\langle W_{R\times T}\rangle_{L,\beta}|\ge\sigma_\eta-(\log C')/R$; let $R\to\infty$ and use Table~\ref{tab:tilted}.
\end{proof}

\begin{remark}[scope]\label{rem:tiltscope}
Like Theorem~\ref{thm:main}, Theorem~\ref{thm:tilted} is proved on tori and for weak limits of torus measures. Free spatial boundaries and
a non-periodic cutting direction are treated neither for Theorem~\ref{thm:main} nor for Theorem~\ref{thm:tilted} (\S\ref{sec:scope}); for
the tilted operator, the compensation at a free boundary that such an extension would need has not been worked out. The weights files
are those of version~1.0 (Proposition~\ref{prop:cert} and Table~\ref{tab:tension}), used without recomputation;
\texttt{c1\_poids\_exacts.py} regenerates them byte for byte.
\end{remark}

\section{The computer-assisted part}\label{sec:cap}

\subsection{Exact enclosures of the weights}\label{sec:weights}
Put $u:=1-\lambda$, so that $M_J(u^2)=\sum_kc_kJ^{2k}u^{2k}$, $c_k=1/(k!(k+1)!)$. Truncate at $K=10$: for $s\in[0,1]$,
$0\le M_J(s)-M_{J,K}(s)\le\tau:=c_{11}J^{22}/(1-J^2/156)$, because $c_{k+1}J^2/c_k=J^2/((k+1)(k+2))\le J^2/156$ for $k\ge11$. Then
\[
\Big(\frac{M_{J,K}(u^2)}{M_{J,K}(1)+\tau}\Big)^3\le h_J(\lambda)\le\Big(\frac{M_{J,K}(u^2)+\tau}{M_{J,K}(1)}\Big)^3 .
\]
The cubes are polynomials in $u$ with rational coefficients, and $\int_0^1(1-u)^{n-1}u\cdot u^j\,du/(n-1)!=(j+1)!/(n+j+1)!$ exactly.
This gives rationals $\At_n^{\inf}\le\At_n\le\At_n^{\sup}$ for $n\le12$ (relative width below $10^{-24}$ at $J_*$). Then
$u_1:=\At_1^{\sup}$, $u_N:=\min(\At_N^{\sup}/\At_{N-1}^{\inf},1/(N+1))$ for $2\le N\le12$, $u_{13}:=1/14$ (which bounds $\rt(N)$ for all
$N\ge13$ by Lemma~\ref{lem:revisit}), $\bar W_n:=\max(u_n,\dots,u_{13})$ and $p_n:=\lceil2^{40}\bar W_n\rceil$. By construction
$\rb_*(N)=p_{\min(N,13)}/2^{40}\ge\sup_{N'\ge N}\rt(N')$, which is item~1 of Proposition~\ref{prop:cert}. Everything is computed with
Python fractions (program \texttt{c1\_poids\_exacts.py}).

\subsection{The integer Perron certificate}\label{sec:perron}
The program \texttt{tm\_cert.c} enumerates one representative per orbit of $\mathfrak G$ on $D^m$ (canonical form, with $s_1$ the most recent step: $s_1=+e_1$;
the first $s_i$ off the $e_1$-axis equal to $+e_2$; the first $s_i$ along the third axis equal to $+e_3$), computes the transitions and window
counts, runs a floating-point power iteration, and rounds the approximate Perron vector to integers $H:=\max(1,\lfloor2^{40}h/
h_{\max}\rfloor)$. It then sets $A:=\max_s\lceil2^{30}\sum_dp_{\min(N(s,d),13)}H(s'_d)/(2^{40}H(s))\rceil$ (exact integer division) and
re-checks \eqref{eq:cert} for every representative in unsigned $128$-bit integer arithmetic; $\rho'=A/2^{30}$. Floating point only
serves to find $H$; the inequality is checked exactly. For $m=11$ there are $7\,821\,056$ representatives and $A=2639854083$.

\emph{Exhaustiveness of the representatives.} The table is closed under transitions (every successor state lies in a represented
orbit), and from any state every state is reached in exactly $m$ transitions (choose the $m$ new steps); so a non-empty closed table
meets every orbit. In addition, its size equals the Burnside count $(6^m+3\cdot4^m+15\cdot2^m)/48$ and no two entries lie in the same
orbit.

\subsection{Independent verifications}\label{sec:indep}
\begin{itemize}
\item \emph{Second program.} \texttt{verif\_cert.c}, written by the same author as \texttt{tm\_cert.c} but with independent
implementation choices, re-verifies the certificate: an orbit is identified by
the minimal base-$6$ code among its $48$ images under explicitly constructed group elements (not by the canonical form of
\texttt{tm\_cert}); window counts are recomputed by explicit tracking of positions; it checks that the entries lie in distinct orbits,
that every transition lands in a represented orbit, that $H\ge1$, and \eqref{eq:cert} in $128$-bit integers. Result for $m=11$: zero
violations, zero missing transitions (also for $m=7,9,13$).
\item \emph{Full state space, $m=11$.} A third program, written in a separate session for an adversarial verification without reusing \texttt{tm\_cert} or
\texttt{verif\_cert}, extends $H$ to all $6^{11}=362\,797\,056$ raw states through orbit look-up and checks \eqref{eq:cert} state by
state, with no symmetry argument: zero violations, largest ratio of the two sides $0.999999999788$. It also recomputes the Perron
root on the raw states: $2.4585557007\ldots<\rho'_*$ \st{numerical}.
\item \emph{Memory $13$.} \texttt{verif\_cert} (zero violations), the third program (all $276\,294\,656$ representatives and
$2\cdot10^8$ random raw states), and a fourth program written by another reviewer (all representatives, $10^7$ random raw states, Burnside count,
weights recomputed in exact rationals with bounded tails): no violation; largest ratio $0.9999999999166$.
\item \emph{Weights.} The bounds $\rb_*(N)\ge\sup_{N'\ge N}\rt(N';J_*)$ were re-derived by a different method (ball arithmetic in Arb
\cite{Joh17}: series of the cube in $(1-\lambda)^2$ with exactly known total sum, $M_J(1)$ from Arb's Bessel functions, and a
rigorous comparison with $p_n/2^{40}$) for $m=11$, $m=13$ and all rows of Table~\ref{tab:tension}; and independent quadratures
(tanh--sinh and Gauss--Legendre) reproduce $p_1,p_2,p_3$ after rounding \st{numerical}.
\item \emph{Arithmetic of item~3.} $J_*\rho'_*<1$ is checked with exact fractions.
\end{itemize}

\subsection{A guard against silent truncation}\label{sec:guard}
In python-flint, operations on power series (\texttt{arb\_series}) are silently truncated at the length \texttt{ctx.cap}, whose
default is $10$, below the $61$ coefficients of the cube of a polynomial of degree $20$. The cross-check \texttt{c2\_poids\_flint.py}
computes the enclosures of \S\ref{sec:weights} in Arb both with polynomials (no truncation) and with series for
$\texttt{cap}\in\{10,30,60,61,80\}$, and requires that the large caps reproduce the polynomial computation and that $\texttt{cap}=10$
differs from it: the relative deviation is $3.8\cdot10^{-7}$ at $\texttt{cap}=10$ (truncation detected) and $0$ at
$\texttt{cap}\ge60$. The Arb balls (200 bits, relative radius below $8\cdot10^{-59}$) have centres within $10^{-14}$ of the exact
rationals of \S\ref{sec:weights}; this is a cross-check, the certification being the exact rational computation. The independent Arb
computation of \S\ref{sec:indep} fixes \texttt{ctx.cap} explicitly and detects the same truncation (a series of length $10$ instead of
$40$).

\subsection{Validation of the counting programs}
With the weights of Borgs and Seiler, $r(1)=\frac12$, $r(N)=1/(N+1)$, the programs reproduce $\rho_1=17/6$, hence $12/17$, and the
values $\rho_3,\dots,\rho_{11}$ of Remark~\ref{rem:contrib}; these values coincide with those of a program written independently
(see the statement on the use of AI), and the numbers of orbits $1,3,11,48,236,1248,6896,39168,226496,1325568,7821056$ ($m=1,\dots,11$) are the Burnside
counts.

\subsection{Tilted certificates (version~1.1)}\label{sec:tiltcap}
The program \texttt{cert\_incline11.c} works on the full state space $D^{m}$ ($6^{11}=362\,797\,056$ states for $m=11$, no symmetry
reduction; base-$6$ code with the most recent step $s_1$ as leading digit). It computes the window counts, runs a floating-point power
iteration for $T_m^{(\eta,e_1)}$ ($50$ to $57$ iterations for the runs of Table~\ref{tab:tilted}), rounds the vector up,
$H:=\max(1,\lceil2^{40}h/h_{\max}\rceil)$, sets $P:=\lceil2^{30}\rho_{\rm fl}(1+10^{-9})\rceil+1$, where $\rho_{\rm fl}$ is the largest ratio
$(T_mh)(s)/h(s)$ of the last iteration, and checks \eqref{eq:tiltcert} for every state in unsigned $128$-bit integers (both sides are
below $2^{107}$ for $a,b\le4096$). For $m=11$ a run takes $34$ to $39$\,s on $8$ cores and about $11$\,GB of memory. With the weights of
Borgs and Seiler and $\eta=0$ it reproduces $\rho_1=17/6$, $\rho_3$ and $\rho_5$ of Remark~\ref{rem:contrib}.

A second program, \texttt{b7\_verif11.py}, re-checks each certificate on all states. It was written separately by the same author
as \texttt{cert\_incline11.c}, but with independent implementation choices (as for \texttt{verif\_cert} in \S\ref{sec:indep}): the
states are indexed with the opposite digit order and $H$ is read through the digit permutation; successors and window counts are recomputed from the
digits, the counts by explicit tracking of positions; and \eqref{eq:tiltcert} is checked exactly with $22$-bit limbs in $64$-bit integers,
without $128$-bit arithmetic. For the five certificates of Table~\ref{tab:tilted} it finds zero violations on all $362\,797\,056$ states;
the smallest relative margins correspond to $4.05$ to $4.56$ units of $P$ (floating-point estimate), and with $P-5$ it reports violations
for each of them; with $P$ lowered by $10^{-5}$ in relative terms it reports violations as well. The numbers $\theta$, $K_0^{(\eta)}$, $C^{(\eta)}$ and $C^{(\eta)}_{\rm odd}$ are computed in exact rationals, and
$\sigma_\eta$ and the logarithms in Arb balls at $128$ and $256$ bits, which are required to overlap (\texttt{b7\_constantes.py}); a
second program written separately during the review of this version (\texttt{checks/table4\_fractions.py}: exact fractions and
$60$-digit decimal logarithms, without Arb) gives the same values.

The six vectors $H$ ($2\,902\,376\,448$ bytes each: the five certificates of Table~\ref{tab:tilted} and the run at $\eta=0$ of
Remark~\ref{rem:twocert}) are supplied, compressed with \texttt{xz}, in the Zenodo record of this version; they are too large for
GitHub. \texttt{cert\_incline11} also regenerates them byte for byte, in about $40$\,s each on $8$ cores, on x86-64 with gcc and no
machine-specific options (checked for all six: the regenerated vectors have the listed fingerprints); their SHA-256 are listed in
\path{supplement_v1.1}. On another platform the floating-point iteration may round differently; a vector obtained in this way is again a
valid certificate if it passes the exact check, but its $H_{\min}$, its $P$ and the constants of Table~\ref{tab:tilted} may then differ
slightly.

\begin{remark}[two certificates at $\eta=0$]\label{rem:twocert}
Run with $\eta=0$ at $J_*$, \texttt{cert\_incline11} returns $P=2\,639\,854\,087$ and $H_{\min}=663\,930\,282\,414$, whereas
Proposition~\ref{prop:cert} has $A=2\,639\,854\,083$ and $H_{\min}=663\,930\,282\,413$. The difference comes from the choice of the integer,
not from the operator or the weights. The floating-point Perron root is $\rho_{\rm fl}=2.45855570066\ldots$, the value found on the raw
states in \S\ref{sec:indep}; thus $2^{30}\rho_{\rm fl}=2\,639\,854\,082.44$, the relative margin of $10^{-9}$ gives
$2\,639\,854\,085.08$, and rounding up and adding $1$ gives $P$. \texttt{tm\_cert} instead takes the smallest integer admissible for its
vector (\S\ref{sec:perron}). With the vector of \texttt{cert\_incline11} the smallest admissible integer is also $A$:
\texttt{b7\_verif11.py} finds zero violations on all $6^{11}$ states with $P=A$, and violations with $P=A-1$. The two vectors come from different floating-point iterations (on orbits in \texttt{tm\_cert}, on raw states here) and
are rounded in opposite directions (down in \texttt{tm\_cert}, up here), hence the difference of one unit in $H_{\min}$. Both certificates are
valid: the published one was re-checked for version~1.1 with \texttt{verif\_cert} (zero violations, ``CERTIFICAT VALIDE'';
\path{supplement_v1.1/data/verif_cert_published_m11.log} records the command and the SHA-256 of the file checked), and the new one
with \texttt{b7\_verif11.py} (zero violations on all $6^{11}$ states). Theorem~\ref{thm:main} uses the published certificate.
\end{remark}

\section{Consistency tests (not part of the proof)}\label{sec:checks}

\subsection{One-dimensional slabs}
On the chain $\Z$, a tree, the ratio of partition functions factorises exactly, $Z(\bs J^\lambda)/Z(J)=\prod_\ell\rho_J(s_\ell)$, so
\eqref{eq:walk} with the exact factor must reproduce the exact correlation $\E[q_0\cdot q_R]=u^R$, $u:=I_2(2J)/I_1(2J)$; these chains
are the slabs of the two-dimensional theory. For two sites this identity was checked analytically (the walks have length $2k+1$ and
the sum reduces to $(J/I_1(2J))\sum_nJ^{2n+1}/(n!(n+2)!)=I_2(2J)/I_1(2J)$). For $R\le4$ and $J\in\{0.2,0.4,0.6\}$ the sum over
\emph{all} walks, organised by edge-crossing numbers without truncation in length, agrees with $u^R$ to $9\cdot10^{-29}$,
$4\cdot10^{-21}$ and $9\cdot10^{-17}$ respectively (truncation of the series in $\lambda$) \st{numerical}. With the Cauchy--Schwarz
factor adapted to two edges per site, $h=\rho_J((1-\lambda)^2)$, the bound $2\Gt(R)$ exceeds $u^R$ by the factors of
Table~\ref{tab:2d}: (I) is not too strong, and its loss is small in this case.

\begin{table}[ht]
\centering\small
\begin{tabular}{@{}lllll@{}}
\toprule
$J\backslash R$ & $1$ & $2$ & $3$ & $4$\\
\midrule
$0.2$ & $1.00443$ & $1.00554$ & $1.00665$ & $1.00776$\\
$0.4$ & $1.01755$ & $1.02196$ & $1.02639$ & $1.03083$\\
$0.6$ & $1.03889$ & $1.04871$ & $1.05862$ & $1.06863$\\
\bottomrule
\end{tabular}
\caption{Ratio $2\Gt(R)/u^R$ on the chain $\Z$, with the site-by-site factor $h=\rho_J((1-\lambda)^2)$ \st{numerical}.}
\label{tab:2d}
\end{table}

\subsection{The \texorpdfstring{$O(4)$}{O(4)} model at order \texorpdfstring{$J^3$}{J\^{}3}: sensitivity of (I)}
For the untwisted $O(4)$ model on $\Z^3$ with weight $\exp(2J\sum q_x\cdot q_y)$, the nearest-neighbour correlation is
$\E[q_0\cdot q_{e_1}]=u+4u^3+O(J^5)=\frac J2+\frac5{12}J^3+O(J^5)$ (the four paths of length three that close a plaquette with the
edge each contribute $u^3$; checked on a ring of four sites by a character expansion). Exact symbolic computation of the walk
expansion at this order gives: with the exact factor $c(\lambda)$, $\frac J2+\frac5{12}J^3$ (equality); with the site-by-site factor
$h=\rho^3$, $\frac J2+\frac{23}{36}J^3$, above the true value; with $h=\rho^k$ the coefficient is $\frac{25}{18}-\frac k4$, below the
true value for $k>35/9$. An error of one whole edge per site in Corollary~\ref{cor:site} (e.g.\ $\rho^4$) would therefore be detected
at this order; an error of half an edge would not.

\subsection{The limit of infinite magnetic coupling}\label{sec:O4}
In the finite-temperature setting of \cite{BS83} (electric and magnetic couplings $J_E,J_M$, $L_0$ time layers), the inequalities of
\S\S\ref{sec:I}--\ref{sec:II} involve only the electric coupling and the positivity of the transfer matrix. For $L_0=1$ and
$J_M\to\infty$ the horizontal links become pure gauge and a slab becomes the untwisted $O(4)$ model on $\Z^3$ with coupling
$2J_E$ per edge ($2J=\beta$ in the isotropic case); the argument then asserts exponential decay of its two-point function for
$2J_E<0.81348$ (or $0.81498$). This is consistent with the Monte Carlo estimates of the critical coupling of the $O(4)$ model on the
simple cubic lattice with action $-\beta\sum s_x\cdot s_y$, $\beta_c=0.935861(8)$ (the Monte Carlo value retained in
\cite[Tables~3--4]{CH98}) and $0.9360(1)$ \cite{KK95}: the margin is a factor $1.15$. This case is not a statement of this note; it is a sanity check,
since no bound that is uniform in $J_M$ can hold beyond the critical coupling of the $O(4)$ model.

\section{Discussion}\label{sec:discussion}

\subsection{Scope and limitations}\label{sec:scope}
\begin{itemize}
\item \emph{Tori only.} Corollary~\ref{cor:site} uses that every site of the spatial torus has exactly six edges, and the trace
formula \eqref{eq:trace} uses periodicity in the cutting direction. For other boundary conditions the factor $h$ must be adapted
(sites with $k<6$ edges receive $\rho^{k/2}\ge\rho^3$), which we have not done. In infinite volume we only obtain the weak limit
points of torus measures; by contrast the range $\beta<2/3$ of \cite[Thm.~7.1]{D7} holds for every Gibbs measure.
\item \emph{Loops and constants.} Only planar rectangular loops are treated, with $R\le L/2$ and $T\le L-1$. The string tension of
Theorem~\ref{thm:main} is tiny near the threshold ($7\cdot10^{-6}$ at $\beta_*$), and the bound is non-trivial only for
$R>R_{\min}\approx1.9\cdot10^6$ there. Theorem~\ref{thm:tilted} (version~1.1) raises the lower bound on the tension at $\beta_*$ to $0.004989$
and lowers $R_{\min}$ to about $3.1\cdot10^{3}$ on every torus ($2.8\cdot10^{3}$ in the weak limits); the method still gives a range rather than a good tension (Remark~\ref{rem:table}).
\item \emph{Tilted bound.} Theorem~\ref{thm:tilted} is proved, like Theorem~\ref{thm:main}, on tori (even and odd) and for weak
limits; free spatial boundaries and a non-periodic cutting direction are not treated (Remark~\ref{rem:tiltscope}). It is not formalised
in Lean; its certificates ($2.9$\,GB each) are supplied in compressed form in the Zenodo record and can be regenerated
(\S\ref{sec:tiltcap}).
\item \emph{Verification.} The proof is computer-assisted and has been checked by machine only (\S\ref{sec:status}). The
certificates for $m=11$ (Proposition~\ref{prop:cert} and the tilted ones) are supplied; the one for $m=13$ is described by its hash.
\item \emph{Limits of the method} \st{numerical}. With the weights of Borgs and Seiler, an exact enumeration of the walks up to length
$14$ and an extrapolation give a growth rate $\mu_\infty\approx2.726$--$2.731$, so that the count alone is limited to about $0.733$;
with the weights of (I), the certified thresholds of Table~\ref{tab:memory} appear to converge slowly to about $0.82$. Memory $15$
would require about $10^{10}$ orbits. Replacing the worst case of six visited neighbours in Corollary~\ref{cor:site} by the actual
number would give $0.8060$ instead of $0.7824$ at memory $1$; this information is not local along the walk, and we have not found a
finite-memory rule that uses it.
\item \emph{Not used.} Borgs and Seiler remark that the perimeter decay of the characters $\Tr U_{C_m}$ could be used to sharpen
their bound on $J$ \cite[second (B.14), p.~377]{BS83}; they give no number, and we do not use it. The improvements
only involve the electric coupling and the positivity of the transfer matrix, so they should apply to the finite-temperature setting
of \cite{BS83} as well; we have not written this out.
\end{itemize}

\subsection{Comparison with heuristic values}\label{sec:heur}
Fr\"ohlich's lecture notes state, ``in mean'' \cite[p.~24, (6.14)--(6.16)]{Fro80a} or ``in measure'' \cite[p.~20,
(34)--(35)]{Fro80b}, a decay estimate $F\lesssim e^{-\alpha|\gamma|}$ for the slab factors of a walk expansion, with
$\alpha>-\ln[(2D-3)(D-1)(4D-4)^{-1}\beta_0]$, and conclude area decay for $\beta_0<4/5$ when $D=4$ \st{read}. \cite[p.~24]{Fro80a}
adds that ``(6.14) is only a rough estimate''; \cite[p.~20]{Fro80b} attributes the estimate to the preprint of \cite{DF80}, whose refined
estimates are stated with proofs announced to ``appear elsewhere'' \cite[preprint pp.~55, 57]{DF80}. We know of no proof, and treat
$4/5$ as a heuristic value.

For $D=4$ the condition reads $5\cdot\beta_0/4<1$. Here $5=2D-3$ is the number of non-backtracking continuations of a walk in
$\Z^3$, and $\beta_0/4$ is the weight of a step to a new site in the walk expansion when $\beta_0$ is the coupling of the slab
$\sigma$-model. In the true slab model this coupling is $\beta$ (Lemma~\ref{lem:slab}), as follows after correcting the factor $2$ of
\cite{DF80} discussed in \cite[\S7.2]{D7}; in this reading the announced range is $\beta<4/5$. Read literally, with $\beta_0$ the
coefficient of $\ReTr$ (\cite[p.~22]{Fro80b} writes $\beta\equiv2\beta_0$), it would be $\beta<8/5$.

In the first reading, $5\beta/4<1$ is the small-$\beta$ linearisation of the following \emph{non-backtracking heuristic}, which we
formulate here (we do not attribute it to a source)
\[
5\,u(\beta)<1,\qquad u(\beta):=\frac{I_2(\beta)}{I_1(\beta)}=\frac\beta4-\frac{\beta^3}{96}+O(\beta^5),
\]
where $u(\beta)$ is the exact correlation $\E[q_x\cdot q_y]$ across a single edge of coupling $\beta$; $5u(\beta)=1$ at
$\beta=0.822225$ \st{numerical}. Our rigorous range $0.81348$ ($0.81498$ with memory $13$) lies above $4/5$ and below $0.8222$
(at $\beta_*$, $5u(\beta_*)=0.9899$), and the thresholds of Table~\ref{tab:memory} seem to approach about $0.82$. The new number is
therefore best read as a rigorous realisation of most of a known heuristic, rather than as an unexpected value.

\section{Reproducibility}\label{sec:repro}
The supplementary folder contains the programs and data below; comments and messages in the programs are in French. SHA-256
fingerprints:

\begin{center}\footnotesize
\begin{tabular}{@{}l@{}}
\toprule
\path{certification/c1_poids_exacts.py}\\
\quad{\scriptsize\texttt{8efd0508f3e21ef0db82df056cfc96b6154d8d14a255590feb4dab7a0b618396}}\\[2pt]
\path{certification/c2_poids_flint.py}\\
\quad{\scriptsize\texttt{6c82f6c30d24a2d993be06fbce6f219f5af4587a80756b94a9db6c1a33c3a939}}\\[2pt]
\path{certification/c3_seuil.py}\\
\quad{\scriptsize\texttt{6c1d5397fa2a5f4dd43b4f0f01437c65d8cf947d22871c988b2dc2af27a0fab1}}\\[2pt]
\path{certification/tm_cert.c}\\
\quad{\scriptsize\texttt{acf257f9344311e74f04540aff7f95cefb55fc8a8e34296e2080cb12fcc98749}}\\[2pt]
\path{certification/verif_cert.c}\\
\quad{\scriptsize\texttt{ebdfacb544f967480c8849da6fad8603f4afa1018dd740793a74a4f2bade33ca}}\\[2pt]
\path{data/poids_m11_J40674_100000.txt}\\
\quad{\scriptsize\texttt{28c789d9a8096a71bc46ba8a9a8bde25e360af4896e0e8f4c6964e9448419317}}\\[2pt]
\path{data/certificat_h_m11_J40674_100000.bin.xz} (8\,780\,360 bytes, \texttt{xz -9e})\\
\quad{\scriptsize\texttt{776f4f27be0aabb32ec81050883448b4004db11289b2a1cf338614e188bbc1e4}}\\[2pt]
\path{data/certificat_h_m11_J40674_100000.bin} (125\,136\,920 bytes, unpacked)\\
\quad{\scriptsize\texttt{b50d3769b0ea17d1da13f4e870a7786407192a9885ecceb390ef18cde50cf142}}\\[2pt]
\path{data/poids_m13_J40749_100000.txt}\\
\quad{\scriptsize\texttt{fcdc6bc36d8fd9de3da60bee15c421ca0a2b0e4119f403dbd39d21085b63c706}}\\[2pt]
\path{certificat_h_m13_J40749_100000.bin} (4\,420\,714\,520 bytes, not supplied)\\
\quad{\scriptsize\texttt{c8ad0deb4a5ec86f8fb3cc5802feaee9a3e1273cb8e14883ccf52ff4590082c7}}\\[2pt]
\bottomrule
\end{tabular}
\end{center}

\medskip
\noindent\emph{The memory-$11$ certificate} is supplied in compressed form, as \path{data/certificat_h_m11_J40674_100000.bin.xz};
running \texttt{xz -dk certificat\_h\_m11\_J40674\_100000.bin.xz} in \path{data/} recreates the file \path{data/certificat_h_m11_J40674_100000.bin}
of the table (same SHA-256) and keeps the compressed file. The uncompressed file is also deposited on Zenodo, in the record of this note.

\noindent\emph{Formats.} A weights file contains $2^{40}$, the number $13$, then $p_1,\dots,p_{13}$. A certificate file contains, in
little-endian 64-bit words, the number $n$ of representatives, their $n$ codes $\sum_id_i6^{m-i}$ (directions $0{:}+e_1$, $1{:}-e_1$,
\dots, $5{:}-e_3$, $d_1$ the most recent step), the $n$ values of $H$, then $A$ and $B=2^{30}$.

\noindent\emph{Commands} (run in \path{supplement/}; \path{w.txt}, \path{c1.json}, \path{H.bin} are output files):
{\footnotesize
\begin{verbatim}
python3 certification/c1_poids_exacts.py 40674 100000 12 10 w.txt c1.json
python3 certification/c2_poids_flint.py 40674 100000 data/c1_J40674.json
gcc -O3 -fopenmp -o certification/tm_cert certification/tm_cert.c -lm
certification/tm_cert 11 w.txt H.bin
gcc -O3 -fopenmp -o certification/verif_cert certification/verif_cert.c
(cd data && xz -dk certificat_h_m11_J40674_100000.bin.xz)
certification/verif_cert 11 data/poids_m11_J40674_100000.txt \
                            data/certificat_h_m11_J40674_100000.bin
\end{verbatim}}
\noindent The first command recreates the weights file byte for byte; the second runs the Arb cross-check with the truncation guard
(it needs a Python interpreter with python-flint installed); \texttt{tm\_cert} recreates a certificate, and \texttt{c3\_seuil.py} runs the dichotomy in $J$ by calling it;
\texttt{xz -dk} unpacks the supplied certificate; the last command prints \texttt{CERTIFICAT VALIDE} (``certificate valid''). Item~3 of Proposition~\ref{prop:cert} is the
Python test
{\footnotesize
\begin{verbatim}
from fractions import Fraction; Fraction(40674,100000)*Fraction(2639854083,2**30) < 1
\end{verbatim}}

\noindent The folder \path{checks/} contains the programs and logs of the independent verifications of \S\S\ref{sec:indep} and
\ref{sec:checks}. The runs used Python~3.14.4, python-flint~0.9.0, gcc~15.2.0 on x86-64. For $m=13$, \texttt{verif\_cert} took
$2$\,min\,$36$\,s on $8$ cores with $10.8$\,GB of resident memory, and the certificate was generated in about $20$ minutes.

\medskip
\noindent\emph{Formalisation in Lean.} The folder \path{lean/} is a Lean~4 project (Lean~4.34.0, Mathlib v4.34.0) that states
Theorem~\ref{thm:main} for the Wilson measure on the torus $(\Z/L\Z)^4$, with the normalised Haar measure on $SU(2)\cong S^3$, and proves
part of its proof; the commands to rebuild it are in \path{supplement/README.md}. The statement is written for the rectangle with a
corner at the origin in the plane of the first and fourth axes (the other loops follow by symmetries of the torus, which are not
formalised), and the expectation is the quotient of the two Haar integrals (the integrability of the integrands, which is clear for bounded
continuous functions on a compact space, is not formalised). Its status is as follows.
\begin{itemize}
\item \emph{Proved in the Lean kernel} (axioms \texttt{propext}, \texttt{Classical.choice} and \texttt{Quot.sound} only; no
\texttt{sorry}, no \texttt{native\_decide}; the main theorems and their closure, $43\,458$ declarations, were also accepted by the
independent kernel nanoda): the combinatorial part of the proof, namely the Perron-type bound on the weighted walk counts,
Lemmas~\ref{lem:window} and~\ref{lem:torus}, the step-by-step reading $\prod_z\At_{n_z}=\prod_j\rt(N_j)$ of \eqref{eq:Gt}, the bound
$|\omega|\ge R$ and the geometric series; the correctness of the certificate checker, for \emph{every} function $H$; the certificates of
memory $m=3$, $4$ and $5$, by \texttt{decide +kernel} on the $6^m$ raw states (those of $m=3,4$ correspond to $\beta\le0.79718$, above
$12/17$, and that of $m=5$ to $\beta\le0.80438$); the final arithmetic in exact rationals ($\theta_*<1$, $3.978<K_0<3.979$,
$563\,935<C_*<563\,936$); the elementary facts $\At_n>0$ and $\At_1\le\frac12$; and the junction theorem
\texttt{theoreme1\_de\_verifiePar}: for every function $H$ such that the checker returns \texttt{true} on the memory-$11$ data, and under
the three hypotheses below, $|\langle W_\gamma\rangle|\le C_*^T\theta_*^{RT}$ for $\beta\le0.81348$.
\item \emph{Established by compiled execution} (the trust is in the Lean and C compilers and in the runtime, not in the kernel): the
memory-$11$ certificate, checked on all $6^{11}=362\,797\,056$ raw states by the very function whose correctness is proved above. The
values of $H$ are read from the file of \texttt{tm\_cert}; since the correctness holds for every $H$, a faulty reading could give
\texttt{true} only if the values read were themselves a valid certificate.
\item \emph{Assumed}, as the three fields of the structure \texttt{HypothesesBS}: (i)~the reduction $|\langle W_\gamma\rangle|\le2^T\Gt^T$
of \eqref{eq:Gt}, which gathers the slab decomposition, the walk representation, the trace formula, Lemma~\ref{lem:pw},
Propositions~\ref{prop:strict} and~\ref{prop:holder} and Corollary~\ref{cor:site} (\S\S\ref{sec:bs}--\ref{sec:I}); (ii)~the monotonicity of
Lemma~\ref{lem:monotone}; (iii)~the certified bounds on the weights, $\rt(N;J_*)\le p_{\min(N,13)}/2^{40}$
(Proposition~\ref{prop:cert}(1); the rational enclosures of \S\ref{sec:weights} are not redone in Lean).
\end{itemize}
Thus the Lean theorem for $\beta\le0.80438$ (memory $5$, the three hypotheses taken at $J_*=40219/10^5$ with the corresponding weights)
rests on these three hypotheses alone, and the theorem for $\beta\le0.81348$ rests on them and on the compiled execution.

\medskip
\noindent\emph{Version~1.1.} The folder \path{supplement_v1.1/} contains the programs of \S\ref{sec:tiltcap}, the weights files at
$J=0.405$, $0.40$ and $0.35294$ (the one at $J_*$ is \path{data/poids_m11_J40674_100000.txt} of \path{supplement/}), the outputs of the
runs, the re-check of the published certificate, and the SHA-256 of the vectors $H$ (\path{data/SHA256_H.txt}); see its
\path{README.md}. SHA-256 fingerprints:

\begin{center}\footnotesize
\begin{tabular}{@{}l@{}}
\toprule
\path{certification/cert_incline11.c}\\
\quad{\scriptsize\texttt{0dc2e7b0c881fee30d44e0bdbc942a0914e1afd4db520205e441b6cc38d3d0c1}}\\[2pt]
\path{certification/b7_verif11.py}\\
\quad{\scriptsize\texttt{38f063a2e49f2e6ef50ca3eb8f860607e818514434a16140e3c2f8d5992f8a07}}\\[2pt]
\path{certification/b7_constantes.py}\\
\quad{\scriptsize\texttt{c494a362d3e426203a3594848d70748076f962c37fd9ce1ab044671c7104f7c8}}\\[2pt]
\path{data/poids_m11_J40500_100000.txt}\\
\quad{\scriptsize\texttt{cc1eda638c4b6cce2c781ab96b067a712f3e9ba5b402480a4a00dd49071fc59c}}\\[2pt]
\path{data/poids_m11_J40000_100000.txt}\\
\quad{\scriptsize\texttt{cf8d8e869249f53512831b8df98c5051aad580912885d6f96bf271cd68ee6b83}}\\[2pt]
\path{data/poids_m11_J35294_100000.txt}\\
\quad{\scriptsize\texttt{17d23b03b618479f3f89cb2eab8757673edcd22454d779e177c1621dc9e0fab4}}\\[2pt]
\path{checks/table4_fractions.py}\\
\quad{\scriptsize\texttt{ccdc7486543c600dc591bd444be8706cbe867d2c4ad87fc6a0a00a9f64854678}}\\[2pt]
\path{data/SHA256_H.txt} (fingerprints of the six vectors $H$, $2\,902\,376\,448$ bytes each, on Zenodo)\\
\quad{\scriptsize\texttt{e0a76870a73d4745d4ae2f1fe02d661d68e081da8ac07f80631a365279436a39}}\\[2pt]
\bottomrule
\end{tabular}
\end{center}

\noindent\emph{Commands} (in \path{supplement_v1.1/}, after copying \path{poids_m11_J40674_100000.txt} from \path{supplement/data/} into
\path{data/}; \path{H.bin} is an output file of $2.9$\,GB):
{\footnotesize
\begin{verbatim}
gcc -O3 -fopenmp -o certification/cert_incline11 certification/cert_incline11.c -lm
certification/cert_incline11 11 data/poids_m11_J40674_100000.txt 201 200 H.bin
python3 certification/b7_verif11.py 11 data/poids_m11_J40674_100000.txt \
        201 200 H.bin 2639868799
sha256sum H.bin
python3 certification/b7_constantes.py data data
python3 checks/table4_fractions.py data
\end{verbatim}}
\noindent The second command prints the certificate of the first row of Table~\ref{tab:tilted} (zero violations), the third re-checks
it, the fourth must give the fingerprint of \path{H_J40674_a201_b200.bin} in \path{data/SHA256_H.txt}, and the last two recompute the
constants of Table~\ref{tab:tilted} with two separate programs (the first of them needs python-flint). The other rows are obtained with
the weights files and the integers $a,b$ of the table. Instead of the second command, the vector can be unpacked from the file
\path{H_J40674_a201_b200.bin.xz} of the Zenodo record (\texttt{xz -dk}), whose name then replaces \path{H.bin} in the next two commands.

\section*{Use of AI}\label{sec:ai}
\addcontentsline{toc}{section}{Use of AI}
The work reported here was carried out with AI models, under the author's direction.
\begin{itemize}
\item \emph{Claude Opus~5.5 (Anthropic)} wrote the proof (\S\S\ref{sec:bs}--\ref{sec:II}), the certification programs
(\texttt{c1}--\texttt{c3}, \texttt{tm\_cert}, \texttt{verif\_cert}) and this note; in a separate session it carried out an
adversarial verification with programs written independently of the certification programs (\S\ref{sec:indep}, \S\ref{sec:checks}).
\item \emph{Kimi K3.1 Pro (Moonshot AI)} wrote, independently, a program for the finite-memory operators with the weights of Borgs
and Seiler, whose values coincide with ours; derived independently the inequality of Proposition~\ref{prop:strict}, with a weaker
site factorisation; carried out an adversarial review of the proof, with recompilation and re-execution of the programs; audited
the memory-$13$ certificate with its own program (\S\ref{sec:indep}); and reviewed a previous version of this note (before the
journal page numbers of \cite{BS83} and the references \cite{Tom07,IS08} were added), regenerating the weights file and the
memory-$11$ certificate byte for byte (no blocking defect; two minor corrections applied); it also reviewed the statement of the Lean
formalisation, the boundary between what is proved, executed and assumed, and the arbitration of the objections (two corrections applied).
\item \emph{DeepSeek V4 Pro (DeepSeek)} contributed exploratory derivations of other approaches and four adversarial reviews of the
proof; their objections led to corrections of the write-up (the orientation of the loop $C_T$ in Proposition~\ref{prop:loop}, the
comparison of exact rather than rounded weights in \S\ref{sec:proof}, and the constant $K_0$). In a last flaw hunt, ten independent
adversarial instances (one per step of the proof) and seven cross-critiques raised eight objections (one of them withdrawn by its
author); the arbitration, by derivation or computation, found no mathematical flaw (the two objections called fatal were refuted by an exact Haar integral and by a counting
argument) and led to three changes of the write-up: the range $0<J\le1$ of Lemma~\ref{lem:pw}, the statement on the string tension
(Corollary~\ref{cor:tension}) and the remark on the range of $R$ (Remark~\ref{rem:table}).
\item \emph{Claude (Anthropic)}, in later sessions, wrote the Lean formalisation of \S\ref{sec:repro} and arbitrated the objections of
the flaw hunt.
\item \emph{Version~1.1.} Claude Opus~5.5 (Anthropic) wrote \S\ref{sec:tilted} and \S\ref{sec:tiltcap}, the three programs of
\S\ref{sec:tiltcap} (\texttt{cert\_incline11.c}, \texttt{b7\_verif11.py} and \texttt{b7\_constantes.py}) and this revision; the use of a tilt and of the
parity for the constants was also proposed, independently, by DeepSeek V4 Pro in an exploratory session. Kimi K3.1 Pro reviewed an
internal write-up of the tilted bound and its programs (recompilation, re-execution at memory~$7$, negative control). The draft of
version~1.1 was reviewed adversarially by DeepSeek V4 Pro (two reviews) and, Kimi being unavailable (quota), by a separate instance
of Claude Opus~5.5 reading the PDF without prior context, which also wrote \texttt{checks/table4\_fractions.py}; see the section on
the changes from version~1.0.
\item \emph{Version~1.2.} The passages on the literature were revised on 11 November 2026 with the assistance of GPT-6.1 Sol (OpenAI), after a review of the draft of 10 November 2026 by Qwen3.8 (Qwen team, run locally); the final text was then read against the sources by Claude (Anthropic) agents on 11 November 2026; no statement, proof, constant or certificate was modified since version~1.1.
\end{itemize}
These reviews are not a substitute for refereeing: no human has checked the proof.

\section*{Changes from version 1.1 (version 1.2)}
\begin{itemize}
\item Literature only: Table~\ref{tab:ranges} gains three rows (Rajabu \cite{Raj26a,Raj26b}) and a note distinguishing the two estimates near $0.730$; \S1.3 gains a paragraph comparing ranges, distinguishing record creation from publication dates, and explicitly making no priority claim; \S1.1 reports the pinning cap of \cite[Prop.~7.1, second preprint]{Raj26b} with its scope. Two bibliography entries are added relative to version~1.1. No statement, proof or constant of our results, certificate, program or file of the supplements is changed.
\item The status of Tomboulis' preprint is no longer described (the words ``unpublished'' and ``has not been published'' are removed); the attribution of the finite-memory walk bound is limited to \cite{PT00}, with \cite{Alm93} cited for related upper bounds; footnote~$^e$ is added to Table~\ref{tab:ranges}; ``(draft of version 1.2)'' is removed from a section title.
\end{itemize}

\section*{Changes from version 1.0}\label{sec:changes}
\addcontentsline{toc}{section}{Changes from version 1.0}
\begin{itemize}
\item New \S\ref{sec:tilted}: the tilted bound (Theorem~\ref{thm:tilted}), its memory-$11$ certificates
(Proposition~\ref{prop:tiltcert}, Table~\ref{tab:tilted}), Corollary~\ref{cor:tilted} and Remark~\ref{rem:tiltscope}. At
$\beta\le0.81348$ the certified lower bound on the string tension rises from $7.05\cdot10^{-6}$ to $0.004989$, and the radius beyond which
the bound is non-trivial on every torus drops from $1.9\cdot10^{6}$ to $3.1\cdot10^{3}$; for $\beta\le0.80$, from $0.0133$ to $0.2253$ and from
$430$ to $36.5$. The threshold, Theorem~\ref{thm:main}, its certificate and the Lean formalisation are unchanged; Theorem~\ref{thm:tilted}
is not formalised.
\item New \S\ref{sec:tiltcap} (programs and verification of the tilted certificates, Remark~\ref{rem:twocert} on the two certificates at
$\eta=0$) and folder \path{supplement_v1.1/} (\S\ref{sec:repro}); the six vectors $H$ of the tilted certificates are deposited,
compressed with \texttt{xz}, in the Zenodo record of this version. The files of version~1.0 are unchanged.
\item Abstract, introduction (\S\S1.2 and~1.4), Remark~\ref{rem:table}, the caption of Table~\ref{tab:tension}, \S\ref{sec:scope} and the
statement on the use of AI updated accordingly. The new Table~\ref{tab:tilted} shifts the former Table~4 to Table~5; section numbers
are unchanged.
\item Review of the draft of version~1.1: two adversarial reviews by DeepSeek V4 Pro (the two objections called blocking were a request
for more explicit wording, now applied, and an arithmetic slip; the other objections were applied or refuted by computation) and a
reading of the PDF without prior context by a separate instance of Claude Opus~5.5 (Kimi K3.1 Pro being unavailable), which found no
mathematical flaw and recomputed Table~\ref{tab:tilted} with its own program; its corrections are included (comparison with
\cite[Cor.~7.2]{D7} in Remark~\ref{rem:table}, rounding in the caption of Table~\ref{tab:tilted}, radii on every torus, scope of
Remark~\ref{rem:tiltscope}, controls at the unit level and Remark~\ref{rem:twocert}). No human has refereed version~1.1.
\end{itemize}


\begin{thebibliography}{99}
\bibitem{Alm93} S.~E.~Alm, \emph{Upper bounds for the connective constant of self-avoiding walks}, Combin.\ Probab.\ Comput.\ 2
(1993) 115--136, doi:10.1017/S0963548300000547.
\bibitem{BS83} C.~Borgs, E.~Seiler, \emph{Lattice Yang--Mills theory at nonzero temperature and the confinement problem}, Comm.\
Math.\ Phys.\ 91 (1983) 329--380, doi:10.1007/BF01208780 (page numbers refer to this version); preprint MPI-PAE/PTh 16/83 (March
1983), with the same text and equation numbers in the passages cited.
\bibitem{CH98} M.~Caselle, M.~Hasenbusch, \emph{The stability of the $O(N)$ invariant fixed point in three dimensions}, J.\ Phys.\ A
31 (1998) 4603--4617, doi:10.1088/0305-4470/31/20/004, arXiv:cond-mat/9711080.
\bibitem{CNS25} S.~Cao, R.~Nissim, S.~Sheffield, \emph{Dynamical approach to area law for lattice Yang--Mills}, arXiv:2509.04688
(2025).
\bibitem{Dou66} R.~G.~Douglas, \emph{On majorization, factorization, and range inclusion of operators on Hilbert space}, Proc.\ Amer.\
Math.\ Soc.\ 17 (1966) 413--415, doi:10.1090/S0002-9939-1966-0203464-1.
\bibitem{DF80} B.~Durhuus, J.~Fr\"ohlich, \emph{A connection between $\nu$-dimensional Yang--Mills theory and $(\nu-1)$-dimensional,
non-linear $\sigma$-models}, Comm.\ Math.\ Phys.\ 75 (1980) 103--151, doi:10.1007/BF01222514; preprint IHES/P/79/302 (1979), to which
``preprint p.'' refers.
\bibitem{FXL18} A.~M.~Ferrenberg, J.~Xu, D.~P.~Landau, \emph{Pushing the limits of Monte Carlo simulations for the three-dimensional
Ising model}, Phys.\ Rev.\ E 97 (2018) 043301, doi:10.1103/PhysRevE.97.043301.
\bibitem{Fro79} J.~Fr\"ohlich, \emph{Confinement in $Z_n$ lattice gauge theories implies confinement in $SU(n)$ lattice Higgs
theories}, Phys.\ Lett.\ B 83 (1979) 195--198, doi:10.1016/0370-2693(79)90684-1.
\bibitem{Fro80a} J.~Fr\"ohlich, \emph{Lectures on Yang--Mills theory}, preprint IHES/P/80/01, Institut des Hautes \'Etudes
Scientifiques (1980).
\bibitem{Fro80b} J.~Fr\"ohlich, \emph{Some comments on the crossover between strong and weak coupling in $SU(2)$ pure Yang--Mills
theory}, Phys.\ Rep.\ 67 (1980) 137--149, doi:10.1016/0370-1573(80)90086-1; preprint IHES/P/80/20 (Les Houches lectures, February
1980), to which page and equation numbers refer.
\bibitem{IS08} K.~R.~Ito, E.~Seiler, \emph{Further discussion of Tomboulis' approach to the confinement problem}, arXiv:0803.3019
(2008).
\bibitem{Joh17} F.~Johansson, \emph{Arb: efficient arbitrary-precision midpoint-radius interval arithmetic}, IEEE Trans.\ Comput.\ 66
(2017) 1281--1292, doi:10.1109/TC.2017.2690633.
\bibitem{KK95} K.~Kanaya, S.~Kaya, \emph{Critical exponents of a three-dimensional $O(4)$ spin model}, Phys.\ Rev.\ D 51 (1995)
2404--2410, doi:10.1103/PhysRevD.51.2404, arXiv:hep-lat/9409001.
\bibitem{OS78} K.~Osterwalder, E.~Seiler, \emph{Gauge field theories on a lattice}, Ann.\ Phys.\ 110 (1978) 440--471,
doi:10.1016/0003-4916(78)90039-8.
\bibitem{PT00} A.~P\"onitz, P.~Tittmann, \emph{Improved upper bounds for self-avoiding walks in $\Z^d$}, Electron.\ J.\ Combin.\ 7
(2000) R21, doi:10.37236/1499.
\bibitem{D7} K.~R\'emondi\`ere, \emph{Dobrushin uniqueness for the four-dimensional $SU(2)$ Wilson lattice gauge measure for
$\beta<2/9$, with interval/Lean certificates}, note v1.5 (1 November 2026; version 1.6 of 12 November 2026 has the same numbering of Section~7), Zenodo, doi:10.5281/zenodo.23856414; Section~7: area law
for $\beta<2/3$.
\bibitem{Raj26a} H.~T.~Rajabu, \emph{Explicit strong-coupling mass gap and a robust landing interface for four-dimensional SU(2) lattice Yang--Mills theory}, Zenodo, doi:10.5281/zenodo.23691339 (27 October 2026; concept doi:10.5281/zenodo.23691338); the archive contains the preprint and companion working notes, among them ``Part~I'' (``An explicit strong-coupling mass gap and area law for four-dimensional SU(2) lattice Yang--Mills: Kantorovich contraction, the sharp single-link constant, and a line gauge'') and ``Part~II'' (``Past the single-site ceiling by exact decimation and the chordal metric: explicit mass gap up to $\beta_W=0.393$ and area law up to $\beta_W=\frac23$ for four-dimensional SU(2) lattice Yang--Mills''), both dated 26 October 2026.
\bibitem{Raj26b} H.~T.~Rajabu, same record title as \cite{Raj26a}, preprint version 2.2 and second preprint ``Multiscale covariance transfer and a conditional continuum mass gap for four-dimensional SU(2) lattice Yang--Mills theory'' (title as printed in that preprint, version 2.1), Zenodo, doi:10.5281/zenodo.23729328 (\texttt{publication\_date} 28 October 2026; record created 29 October 2026, 00:48 UTC; a new version of \cite{Raj26a}).
\bibitem{Tom07} E.~T.~Tomboulis, \emph{Confinement for all values of the coupling in four-dimensional $SU(2)$ gauge theory},
arXiv:0707.2179 (2007).
\bibitem{Wil74} K.~G.~Wilson, \emph{Confinement of quarks}, Phys.\ Rev.\ D 10 (1974) 2445--2459, doi:10.1103/PhysRevD.10.2445.
\end{thebibliography}

\end{document}
